\exercisesection{}

\begin{enumerate}
\def\labelenumi{\arabic{enumi})}
\item
  Let G be a locally compact group. Let \(\pi\) be a unitary representation of G in a complex Hilbert space E and \(x \in \mathrm{E}\) . We have \(\pi(g)x = x\) for all \(g \in \mathrm{G}\) if and only if the positive-definite function \(g \mapsto \langle x | \pi(g)x \rangle\) is constant.
\item
  Let G be a locally compact group. For every positive-definite function \(\varphi\) on G, we denote by \((\pi_{\varphi}, x_{\varphi})\) a cyclic Hilbert representation of \(\varphi\) .
\end{enumerate}

\begin{enumerate}
\def\labelenumi{\alph{enumi})}
\item
  Let \(\varphi_1, \varphi_2\) and \(\varphi_3\) be positive-definite functions on G. If \(\varphi_3 = \varphi_1 + \varphi_2\) , then there exists an injective G-morphism from \(\pi_{\varphi_1}\) into \(\pi_{\varphi_3}\) .
\item
  If \(\varphi\) is a constant nonzero positive-definite function, then \(\pi_{\varphi}\) is trivial of dimension 1.
\item
  For every unitary representation \(\pi\) of G in a Hilbert space E, and for every \(x \in \mathrm{E}\) , there exists an injective G-morphism from \(\pi_{\varphi}\) into \(\pi\) , where \(\varphi(g) = \langle x \mid \pi(g)x \rangle\) .
\item
  Let H be a closed subgroup of G and let \(\varphi\) be a positive-definite function on G. Denote by \(\psi\) the restriction of \(\varphi\) to H; this is a positive-definite function on H, and there exists an injective H-morphism from \(\pi_{\psi}\) into \(\operatorname{Res}_{\mathrm{H}}^{\mathrm{H}}(\pi)\) .
\end{enumerate}

\begin{enumerate}
\def\labelenumi{\arabic{enumi})}
\setcounter{enumi}{2}
\tightlist
\item
  \begin{enumerate}
  \def\labelenumii{\alph{enumii})}
  \tightlist
  \item
    Give an example of an involutive Banach algebra admitting a positive linear form which is not continuous.
  \end{enumerate}
\end{enumerate}

\begin{enumerate}
\def\labelenumi{\alph{enumi})}
\setcounter{enumi}{1}
\tightlist
\item
  Give an example of an involutive Banach algebra A, a continuous positive linear form \(\lambda\) on A, and Hilbert realizations \((\mathrm{E}_{1}, e_{1}, \varphi_{1})\) and \((\mathrm{E}_{2}, e_{2}, \varphi_{2})\) of \(\lambda\) such that \((\mathrm{E}_{1}, e_{1}, \varphi_{1})\) is a cyclic Hilbert realization and such that the unique linear map u of prop. 6 of V, p.~440 does not satisfy \(u(e_{1}) = e_{2}\) .
\end{enumerate}

\begin{enumerate}
\def\labelenumi{\arabic{enumi})}
\setcounter{enumi}{3}
\tightlist
\item
  Let A and B be star algebras. A linear map u from A to B is said to be positive if \(u(\mathrm{A}_{+}) \subset \mathrm{B}_{+}\) , and completely positive if, for every integer \(k \in N\) , the map \(u_{k} = 1_{C^{k}} \otimes u: C^{k} \otimes A \to C^{k} \otimes B\) is positive.
\end{enumerate}

\begin{enumerate}
\def\labelenumi{\alph{enumi})}
\tightlist
\item
  Let A be a unital star algebra, E and F complex Hilbert spaces, v a linear map from E to F, and \(\varphi\) a unital morphism of involutive algebras from A to \(\mathcal{L}(\mathrm{F})\) such that \(\|\varphi\|^{2} = \|\varphi(1)\|\) . The map \(u: A \to \mathcal{L}(E)\) defined by \(u(a) = v^{*} \circ \varphi(a) \circ v\) is completely positive.
\end{enumerate}

The goal of the exercise is to prove the converse of this assertion (“Stinespring theorem”). Let A be a unital star algebra, E a complex Hilbert space, and \(u \colon \mathrm{A} \to \mathcal{L}(\mathrm{E})\) a completely positive linear map.

\begin{enumerate}
\def\labelenumi{\alph{enumi})}
\setcounter{enumi}{1}
\tightlist
\item
  Let \(F = A \otimes E\) . There exists a positive sesquilinear form q on F such that
\end{enumerate}

\[
q (a \otimes x, b \otimes y) = \langle x | \varphi (a ^ {*} b) y \rangle
\]

for all \((a,b)\in \mathrm{A}^2\) and \((x,y)\in \mathrm{E}^2\) .

Let \(\widehat{\mathbf{F}}\) be the separated completion Hilbert space of \(\mathbf{F}\) for the sesquilinear form \(q\) .

\begin{enumerate}
\def\labelenumi{\alph{enumi})}
\setcounter{enumi}{2}
\tightlist
\item
  Let \(b \in A\) . The endomorphism of F such that \(a \otimes x \mapsto ba \otimes x\) for all \(a \in A\) and \(x \in E\) induces by passage to quotients an endomorphism \(\varphi(b)\) of \(\widehat{F}\) such that \(\|\varphi(b)\|\leqslant\|b\|\) ; the map \(\varphi\colon A\to\mathcal{L}(F)\) is a unital morphism of involutive algebras.
\end{enumerate}

\begin{enumerate}
\def\labelenumi{\alph{enumi})}
\setcounter{enumi}{3}
\tightlist
\item
  Let \(v\) be the linear map from E to \(\widehat{\mathrm{F}}\) such that \(v(x)\) is the class of \(1 \otimes x\) for all \(x \in \mathrm{E}\) . We have \(\|v\|^2 = \| \varphi(1)\|\) , and \(u(a) = v^* \circ \varphi(a) \circ v\) for all \(a \in \mathrm{A}\) .
\end{enumerate}

\begin{enumerate}
\def\labelenumi{\arabic{enumi})}
\setcounter{enumi}{4}
\tightlist
\item
  Let X be a separated topological space. A continuous function \(f: X \times X \to R\) is said to be of conditionally negative type if it satisfies the conditions
\end{enumerate}

\begin{enumerate}
\def\labelenumi{(\roman{enumi})}
\item
  For every \(x \in X\) , we have \(f(x, x) = 0\) ;
\item
  For all \(x\) and \(y\) in \(X\) , we have \(f(x,y) = f(y,x)\) ;
\item
  For every integer \(n \in N\) , every family \((x_{i})_{0 \leqslant i \leqslant n}\) of elements of X and every family \((t_{i})_{0 \leqslant i \leqslant n}\) of real numbers such that
\end{enumerate}

\[
\sum_ {i = 1} ^ {n} t _ {i} = 0,
\]

we have

\[
\sum_ {i = 0} ^ {n} \sum_ {j = 0} ^ {n} t _ {i} t _ {j} f (x _ {i}, x _ {j}) \leqslant 0.
\]

We denote by Noy \(_{-}\) (X) the set of functions of conditionally negative type on X × X.

\begin{enumerate}
\def\labelenumi{\alph{enumi})}
\item
  Let E be a real Hilbert space and \(\varphi\colon\mathrm{X}\to\mathrm{E}\) a continuous map. The function \(f(x,y)=\|\varphi(x)-\varphi(y)\|^{2}\) is of conditionally negative type.
\item
  Let \(f\in\mathrm{Noy}_{-}(\mathrm{X})\) and \(x_{0}\in\mathrm{X}\) . There exists a real Hilbert space E and a continuous map \(\varphi\colon\mathrm{X}\to\mathrm{E}\) such that
  \begin{enumerate}
  \def\labelenumii{\alph{enumii})}
  \tightlist
  \item
    For every \((x,y)\in \mathrm{X}\times \mathrm{X}\) , we have \(f(x,y) = \| \varphi (x) - \varphi (y)\| ^2\) ;
  \item
    The family \((\varphi(x)-\varphi(x_{0}))_{x\in\mathrm{X}}\) is total in E.
  \end{enumerate}
\end{enumerate}

We say that \((\mathrm{E},\varphi)\) is a cyclic Hilbert representation of f.

\begin{enumerate}
\def\labelenumi{\alph{enumi})}
\setcounter{enumi}{2}
\item
  Let \(f \in \mathrm{Noy}_{-}(\mathrm{X})\) and \(x_0 \in \mathrm{X}\) . Let \((\mathrm{E}, \varphi)\) and \((\mathrm{E}', \varphi')\) be cyclic Hilbert representations of \(f\) . There exists a unique affine isometry \(u: \mathrm{E} \to \mathrm{E}'\) such that \(\varphi' = u \circ \varphi\) .
\item
  The set \(\mathrm{Noy}_{-}(\mathrm{X})\) is a convex cone in \(\mathcal{C}(\mathrm{X}\times \mathrm{X};\mathbf{R})\) . For every universally positive kernel \(g\in \mathrm{Noy}_{+}(\mathrm{X})\) such that \(g\) is real-valued, the map \(f(x,y) = g(x,x) - g(x,y)\) belongs to \(\mathrm{Noy}_{-}(\mathrm{X})\) .
\item
  Let \(\mathfrak{F}\) be a filter on \(\mathrm{Noy}_{-}(\mathrm{X})\) which converges simply to \(f\colon \mathrm{X}\times \mathrm{X}\to \mathbf{R}\) . If \(f\) is continuous, then \(f\in \mathrm{Noy}_{-}(\mathrm{X})\) .
\item
  Let \(f: X \times X \to R\) be a continuous function satisfying conditions (i) and (ii) above. Let \(x_{0} \in X\). Let \(g: X \times X \to R\) be such that \(g(x, y) = f(x, x_{0}) + f(y, x_{0}) - f(x, y)\) . Then \(f \in \text{Noy}_{-}(X)\) if and only if \(g \in \text{Noy}_{+}(X)\) .
\item
  Suppose X is nonempty. Let \(f: \mathrm{X} \times \mathrm{X} \to \mathbf{R}\) be a continuous function satisfying conditions (i) and (ii) above. Then \(f \in \mathrm{Noy}_{-}(\mathrm{X})\) if and only if we have \(\exp(-tf) \in \mathrm{Noy}_{+}(\mathrm{X})\) for all \(t \in \mathbf{R}_{+}\) (“Schoenberg theorem”). (To prove that \(\exp(-f)\) is of positive type when \(f \in \mathrm{Noy}_{-}(\mathrm{X})\) , fix \(x_0 \in \mathrm{X}\) ; show that
\end{enumerate}

\[
g (x, y) = \exp (f (x, x _ {0}) + f (y, x _ {0}) - f (x, y))
\]

and

\[
h (x, y) = \exp (- f (x, x _ {0})) \exp (- f (y, x _ {0}))
\]

belong to Noy \(_{+}\) (X).)

If G is a topological group, we say that a function \(\varphi \in \mathcal{C}(\mathrm{G})\) is of negative type if the function \(f\) on \(\mathrm{G} \times \mathrm{G}\) defined by \(f(g,h) = \varphi(h^{-1}g)\) is a conditionally negative-definite kernel.

\begin{enumerate}
\def\labelenumi{\arabic{enumi})}
\setcounter{enumi}{5}
\tightlist
\item
  Let G be a discrete topological group. Let A be the involutive subalgebra of \(\mathcal{L}(\ell^{2}(\mathrm{G}))\) generated by the image of the left regular representation \(\lambda\) . This is a unital star algebra.
\end{enumerate}

For every function of positive type \(\varphi\in\mathrm{Pos}(\mathrm{G})\) , there exists a unique completely positive linear map \(u_{\varphi}\) from A to A (cf.~exercise 4) such that \(u_{\varphi}(\boldsymbol{\lambda}(g))=\varphi(g)\boldsymbol{\lambda}(g)\) for all \(g\in\mathrm{G}\) (Use a cyclic Hilbertian representation of \(\varphi\) .)

\begin{enumerate}
\def\labelenumi{\arabic{enumi})}
\setcounter{enumi}{6}
\tightlist
\item
  Let I be a finite set and \(G = F(I)\) the free group constructed on I (A, I, p.~84). Let A be the involutive subalgebra of \(\mathcal{L}(\ell^{2}(G))\) generated by the image of the left regular representation \(\lambda\) .
\end{enumerate}

\begin{enumerate}
\def\labelenumi{\alph{enumi})}
\item
  For every \(g \in G\) , we denote by \(|g|\) the length of the unique reduced decomposition of g (A, I, p.~146, exercise 19). The map \(d(g, h) = |h^{-1}g|\) is a metric on G.
\item
  Let \(\lambda \in \mathbf{R}_{+}^{*}\) . The function on G defined by \(\psi_{\lambda}(g) = \exp(-\lambda|g|)\) is a positive-definite function. (Show that the function \(g \mapsto |g|\) is a function of negative type, cf.~exercise 5.)
\item
  For every integer \(k \in N\) , let \(G_{k}\) be the set of \(g \in G\) such that \(|g| = k\) and let \(\varphi_{k}\) be the characteristic function of \(G_{k}\) . Let f and g be functions on G with support in \(G_{k}\) and \(G_{l}\) respectively. For every \(m \in N\) , we have \((f * g)\varphi_{m} = 0\) unless \(k + l - m\) is even and \(|k - l| \leqslant m \leqslant k + l\) ; in this case, we have
\end{enumerate}

\[
\left\| (f * g) \varphi_ {m} \right\| \leqslant \| f \| \| g \|.
\]

\begin{enumerate}
\def\labelenumi{\alph{enumi})}
\setcounter{enumi}{3}
\tightlist
\item
  Let f be a function on G with finite support. We have
\end{enumerate}

\[
\| \boldsymbol {\lambda} (f) \| \leqslant \sum_ {n \in \mathbf {N}} (n + 1) \| f \varphi_ {n} \|,
\]

and

\[
\left\| \boldsymbol {\lambda} (f) \right\| \leqslant 2 \left(\sum_ {g \in G} (1 + | g |) ^ {4} | f (g) | ^ {2}\right) ^ {1 / 2}.
\]

\begin{enumerate}
\def\labelenumi{\alph{enumi})}
\setcounter{enumi}{4}
\item
  Let \(\psi\) be a function on G such that \(M = \sup_{g \in G} |\psi(g)| (1 + |g|)^2\) is finite. There exists a unique linear map \(u_\psi\) from A to A such that \(u_\psi(\boldsymbol{\lambda}(g)) = \psi(g)\boldsymbol{\lambda}(g)\) for all \(g \in G\) . We have \(\|u_\psi\| \leqslant 2M\) .
\item
  Let X be the set of functions \(\psi\) on G with finite support such that \(\|u_{\psi}\|\leqslant1\) . There exists a filter \(\mathfrak{F}\) on X such that
\end{enumerate}

\[
\lim _ {\psi , \mathfrak {F}} u _ {\psi} = 1 _ {\mathrm{A}}
\]

in the space \(\mathcal{L}(\mathrm{A})\) endowed with the topology of simple convergence. In particular, the Banach space underlying A has the approximation property. (Consider the functions \(\psi_{\lambda,n}(g) = (\varphi_0 + \dots + \varphi_n)\psi_\lambda\) for \(\lambda \in \mathbf{R}_+^*\) and \(n \in \mathbf{N}\) , and use the preceding exercise.)

\begin{enumerate}
\def\labelenumi{\arabic{enumi})}
\setcounter{enumi}{7}
\tightlist
\item
  One identifies the group \(\mathbf{R}\) and its dual group by the map \((x,y)\mapsto\) \(\exp (i\pi xy)\) . Let \(\alpha \in \mathbf{R}_{+}\) .
\end{enumerate}

\begin{enumerate}
\def\labelenumi{\alph{enumi})}
\item
  There exists a positive measure on \(\mathbf{R}\) whose Fourier transform is the function \(f_{\alpha}(x) = \exp(-|x|^{\alpha})\) if and only if \(\alpha \in ]0,2]\) . When this is the case, one denotes by \(\mu_{\alpha}\) this measure, which is unique; it is called the stable measure of parameter \(\alpha\) .
\item
  The measure \(\mu_{\alpha}\) has total mass 1; the identity function f on R is integrable with respect to \(\mu_{\alpha}\) if and only if \(\alpha = 2\) .
\end{enumerate}

\begin{enumerate}
\def\labelenumi{\arabic{enumi})}
\setcounter{enumi}{8}
\tightlist
\item
  The goal of this exercise is to present another proof of Theorem 5 of V, p.~455. Let G be a locally compact commutative topological group, equipped with a Haar measure \(\mu\) . We denote by \(\widehat{f}\) the Fourier transform of \(f \in \mathrm{L}^1(\mathrm{G})\) , and we denote by \(\mathcal{F}(\mathrm{G}) \subset \mathcal{C}_0(\widehat{\mathrm{G}})\) the image of the Fourier transform on \(\mathrm{L}^1(\mathrm{G})\) .
\end{enumerate}

\begin{enumerate}
\def\labelenumi{\alph{enumi})}
\item
  The space \(\mathcal{F}(\mathrm{G})\) is a dense subspace of \(\mathcal{C}_0(\mathrm{G})\) .
\item
  For every \(f \in \mathrm{L}^{1}(\mathrm{G})\) , the spectral radius of f in the Banach algebra \(\mathrm{L}^{1}(\mathrm{G})\) is equal to \(\|\widehat{f}\|_{\infty}\) .
\item
  Let \(\varphi\in\mathrm{Pos}(\mathrm{G})\) . It is a uniformly continuous and bounded function on \(\mathrm{G}\) ; for every \(f\in\mathrm{L}^{1}(\mathrm{G})\) , we have
\end{enumerate}

\[
\int_ {G} \int_ {G} f (x) \overline {{f (y)}} \varphi (x - y) d (\mu \otimes \mu) (x, y) \geqslant 0.
\]

\begin{enumerate}
\def\labelenumi{\alph{enumi})}
\setcounter{enumi}{3}
\tightlist
\item
  Let \(\lambda\) be the continuous linear form on \(\mathrm{L}^{1}(\mathrm{G})\) determined by \(\varphi\) and we set \(b(f,g)=\lambda(f^{*}*g)\) for all \((f,g)\in\mathrm{L}^{1}(\mathrm{G})\times\mathrm{L}^{1}(\mathrm{G})\) . The map \(b\) is a positive hermitian form on \(\mathrm{L}^{1}(\mathrm{G})\) .
\end{enumerate}

\begin{enumerate}
\def\labelenumi{\alph{enumi})}
\setcounter{enumi}{4}
\tightlist
\item
  For every \(f \in \mathrm{L}^1(\mathrm{G})\) , we have \(|\lambda(f)|^2 \leqslant b(f, f)\) .
\end{enumerate}

\begin{enumerate}
\def\labelenumi{\alph{enumi})}
\setcounter{enumi}{5}
\tightlist
\item
  Let \(f \in \mathrm{L}^{1}(\mathrm{G})\) and \(g = f^{*} * f \in \mathcal{C}_{b}(\mathrm{G})\) . One defines \(g_{1} = g\) and \(g_{n+1} = g_{n} * g\) for every integer \(n \geqslant 1\) . For every integer \(n \geqslant 1\) , we have \(|\lambda(f)| \leqslant \|g_{n}\|_{1}^{2^{-n}}\) .
\end{enumerate}

\begin{enumerate}
\def\labelenumi{\alph{enumi})}
\setcounter{enumi}{6}
\tightlist
\item
  For every \(f \in \mathrm{L}^{1}(\mathrm{G})\) , we have \(|\lambda(f)| \leqslant \| \widehat{f} \|_{\infty}\) . There exists a continuous linear map \(\widehat{\lambda}\) on \(\mathcal{F}(\mathrm{G})\) such that \(\lambda(f) = \widehat{\lambda}(\widehat{f})\) for all \(f \in \mathrm{L}^{1}(\mathrm{G})\) .
\end{enumerate}

\begin{enumerate}
\def\labelenumi{\alph{enumi})}
\setcounter{enumi}{7}
\tightlist
\item
  There exists a bounded measure \(\nu\) on \(\widehat{G}\) such that \(\lambda(f) = \nu(\widehat{f})\) for every \(f \in \mathrm{L}^{1}(\mathrm{G})\) ; the measure \(\nu\) is positive and one has
\end{enumerate}

\[
\varphi (x) = \int_ {\widehat {\mathrm{G}}} \langle \chi , x \rangle d \nu (\chi)
\]

for all \(x \in G\) .

\begin{enumerate}
\def\labelenumi{\arabic{enumi})}
\setcounter{enumi}{9}
\tightlist
\item
  \begin{enumerate}
  \def\labelenumii{\alph{enumii})}
  \tightlist
  \item
    Let \(\mathrm{G} = \mathbf{R}\) . The set \(\operatorname{Pos}_1(\mathrm{G})\) is not compact for the weak topology.
  \end{enumerate}
\end{enumerate}

\begin{enumerate}
\def\labelenumi{\alph{enumi})}
\setcounter{enumi}{1}
\tightlist
\item
  Let G = R/Z. The weak topology on \(\operatorname{Pos}_{\leqslant1}(\mathbf{G})\) does not coincide with the topology of compact convergence.
\end{enumerate}

\begin{enumerate}
\def\labelenumi{\arabic{enumi})}
\setcounter{enumi}{10}
\tightlist
\item
  Let G be a locally compact group.
\end{enumerate}

Let \(\varrho\) and \(\pi\) be unitary representations of G. One says that \(\varrho\) is weakly contained in \(\pi\) , and one writes \(\varrho \leqslant \pi\) , if the set of diagonal matrix coefficients of \(\varrho\) is contained in the closure of the set of finite sums of diagonal matrix coefficients of \(\pi\) in the space \(\mathcal{C}(\mathrm{G})\) endowed with the topology of compact convergence.

\begin{enumerate}
\def\labelenumi{\alph{enumi})}
\item
  The relation « \(\varrho\) and \(\pi\) are unitary representations of G and \(\varrho \leqslant \pi\) » is a preorder relation. If \(\varrho'\) (resp. \(\pi'\) ) is isomorphic to \(\varrho\) (resp. \(\pi\) ) then \(\varrho \leqslant \pi\) if and only if \(\varrho' \leqslant \pi'\) .
\item
  If \(\varrho\) is isomorphic to a subrepresentation of \(\pi\) , then \(\varrho \leqslant \pi\) . If I and J are nonempty finite sets, and if \(\varrho \leqslant \pi\) , then \(\bigoplus_{i \in I} \varrho \leqslant \bigoplus_{j \in J} \pi\) .
\item
  Let X be a subset of the space E of \(\varrho\) which generates \(\pi\) . We have \(\varrho \leqslant \pi\) if and only if every matrix coefficient \(g \mapsto \langle x | \varrho(g)x \rangle\) with \(x \in X\) belongs to the closure of the set of finite sums of diagonal matrix coefficients of \(\pi\) in the space \(\mathcal{C}(\mathrm{G})\) endowed with the topology of compact convergence. (Show that the set Y of elements of E satisfying this property is a closed vector subspace of E.)
\item
  Let H be a closed subgroup of G. Let \(\pi\) and \(\varrho\) be unitary representations of G such that \(\varrho \leqslant \pi\) ; we then have \(\operatorname{Res}_{\mathrm{H}}^{\mathrm{G}}(\varrho) \leqslant \operatorname{Res}_{\mathrm{H}}^{\mathrm{G}}(\pi)\) .
\item
  Let H be a closed subgroup of G and let \(\nu\) be a Haar measure on H. Let \(\pi\) and \(\varrho\) be unitary representations of H such that \(\varrho \leqslant \pi\) ; we then have \(\operatorname{Ind}_{\mathrm{H}}^{\mathrm{G}}(\varrho) \leqslant \operatorname{Ind}_{\mathrm{H}}^{\mathrm{G}}(\pi)\) . (Let \(\varrho' = \operatorname{Ind}_{\mathrm{H}}^{\mathrm{G}}(\varrho)\) and \(\pi' = \operatorname{Ind}_{\mathrm{H}}^{\mathrm{G}}(\pi)\) ; using part (c), reduce to approximating matrix coefficients \(g \mapsto \langle f | \varrho'(g)f \rangle\) for f of the form
\end{enumerate}

\[
f (x) = \int_ {\mathrm{H}} f (x h) \varrho (h) y d \nu (y)
\]

where \(\varphi\in\mathcal{K}(\mathrm{G})\) and \(y\) belongs to the space of \(\varrho\) .)

\begin{enumerate}
\def\labelenumi{\arabic{enumi})}
\setcounter{enumi}{11}
\item
  Let G = R and let \(\pi\) be the regular representation of G in \(\mathrm{L}^{2}(\mathbf{R})\) . Prove that \(\chi \leqslant \pi\) for every character \(\chi \in \widehat{G}\) .
\item
  Let \(\varrho\) be an irreducible unitary representation of G in a Hilbert space E, and \(\pi\) a unitary representation of G in F. Let \(x \in E\) be of norm 1. One assumes that \(\varrho \leqslant \pi\) (cf.~Exercise 11).
\end{enumerate}

We denote by \(A \subset \mathcal{C}(G)\) the set of normalized diagonal matrix coefficients of \(\pi\) , and \(\widetilde{A}\) the closed convex hull of A in \(L^{\infty}(G)\) for the weak topology.

Let \(x \in E\) be of norm 1 and let \(f: g \mapsto \langle x \mid \varrho(g)x \rangle\) be the associated diagonal matrix coefficient.

\begin{enumerate}
\def\labelenumi{\alph{enumi})}
\item
  The set \(\widetilde{\mathrm{A}}\) is a compact convex set in \(\mathrm{L}^{\infty}(\mathrm{G})\) and \(f\) is an extreme point of \(\widetilde{\mathrm{A}}\) .
\item
  The function f belongs to the closure of A in \(\widetilde{A}\) for the weak topology.
\item
  The function f belongs to the closure of A in \(\mathcal{C}(G)\) endowed with the topology of compact convergence.
\item
  The trivial representation of G in C is weakly contained in \(\pi\) if and only if, for every compact subset C of G and every real \(\varepsilon > 0\) , there exist \(x \in F\) of norm 1 such that
\end{enumerate}

\[
\sup _ {g \in \mathrm{C}} \| \pi (g) x - x \| <   \varepsilon .
\]

\begin{enumerate}
\def\labelenumi{\arabic{enumi})}
\setcounter{enumi}{13}
\tightlist
\item
  Let G be a locally compact group. One says that G has the property ]T[ of Kazhdan \(^{(1)}\) if a unitary representation \(\pi\) of G has a nonzero trivial subrepresentation \(^{(2)}\) if and only if the trivial representation \(\mathbf{1}_{\text{G}}\) of G on C is weakly contained in \(\pi\) .
\end{enumerate}

\begin{enumerate}
\def\labelenumi{\alph{enumi})}
\item
  Let G be a locally compact group having property ]T[. Let H be the set of open subgroups of G that are generated by a compact subset of G. The group G is the union of the subgroups H ∈ H.
\item
  Let \(\pi\) be the representation of G that is the Hilbert sum of the quasi-regular representations of G on G/H for \(H \in H\) (cf. exercise 2 of V, p. 486). We have \(1_{G} \leqslant \pi\) .
\item
  The group G is generated by a compact subset of G. In particular, if G is discrete, then G is generated by a finite subset of G.
\end{enumerate}

\begin{enumerate}
\def\labelenumi{\arabic{enumi})}
\setcounter{enumi}{14}
\tightlist
\item
  Let G be a locally compact group, C a subset of G, and \(\varepsilon > 0\) . One says that \((\mathrm{C}, \varepsilon)\) is a Kazhdan pair for G if, for every unitary representation \(\pi\) of G in a Hilbert space E, the existence of a vector \(x \in E\)
\end{enumerate}

such that

\[
\sup _ {g \in \mathrm{C}} \| \pi (g) x - x \| <   \varepsilon \| x \|,
\]

  implies the existence of a nonzero trivial subrepresentation of \(\pi\) . We say that C is a Kazhdan set if there exists \(\varepsilon > 0\) such that \((C, \varepsilon)\) is a Kazhdan pair.

\begin{enumerate}
\def\labelenumi{\alph{enumi})}
\tightlist
\item
  Let \((\mathrm{C}, \varepsilon)\) be a Kazhdan pair of G. Let \(\pi\) be a unitary representation of G in a Hilbert space E. Let \(x \in E\) be nonzero and \(\delta > 0\) such that
\end{enumerate}

\[
\sup _ {g \in \mathrm{C}} \| \pi (g) x - x \| <   \delta \varepsilon \| x \|.
\]

We then have \(\|p(x)-x\|\leqslant\delta\|x\|\) , where p is the orthogonal projector of E whose image is the space of G-invariant vectors.

\begin{enumerate}
\def\labelenumi{\alph{enumi})}
\setcounter{enumi}{1}
\item
  A locally compact group G has property ]T[ if and only if there exists a Kazhdan pair (C,ε). (To prove that property ]T[ implies the assertion, proceed by contraposition by considering the unitary representation that is the Hilbert sum of representations π \(_{C,\varepsilon}\) parametrized by a compact subset C of G and ε > 0 such that π \(_{C,\varepsilon}\) contains no nontrivial subrepresentation but there exists x \(_{C,\varepsilon}\) of norm 1 with sup \(_{g\in\mathbf{C}}\|\pi(g)x_{\mathrm{C},\varepsilon}-x_{\mathrm{C},\varepsilon}\|<\varepsilon.\)
\item
  Let G be a locally compact group having property ]T[. Every compact subset C of G that generates G is a Kazhdan set.
\item
  Let G be a locally compact group having property ]T[ and let C be a compact subset of G with nonempty interior that is a Kazhdan set. Then C generates G. In particular, if G is discrete, then the Kazhdan sets are precisely the finite generating subsets of G. (The subgroup H of G generated by C is open; consider the quasi-regular representation \(\pi\) of G on G/H and show that the projection of the characteristic function \(\varphi\) of H \(\in\) G/H onto the space of invariant vectors of \(\pi\) is equal to \(\varphi\) .)
\end{enumerate}

\begin{enumerate}
\def\labelenumi{\arabic{enumi})}
\setcounter{enumi}{15}
\item
  Let G be a locally compact group that is countable at infinity and let H be a closed subgroup of G. If G has property ]T[, and if G/H admits a bounded G-invariant measure, then H has property ]T[. (Use Exercise 20 of V, p.~492.)
\item
  Let \(n \geqslant 2\) be an integer. We consider the real Lie group \(\mathbf{G} = \mathbf{R}^n \rtimes \mathbf{SL}(n, \mathbf{R})\) , and we denote by H the subgroup \(\mathbf{R}^n \times \{e\}\) of \(\mathbf{G}\) , identified with \(\mathbf{R}^n\) .
\end{enumerate}

Let \(\pi\) be a unitary representation of G in a complex Hilbert space E. Suppose that for every compact subset C of G and every \(\varepsilon > 0\) , there exists \(x \in E\) nonzero such that

\[
\sup _ {g \in \mathrm{C}} \| \pi (g) x - x \| <   \varepsilon \| x \|.
\]

\begin{enumerate}
\def\labelenumi{\alph{enumi})}
\item
  There exists a sequence \((x_{m})\) of vectors of norm 1 in E such that the functions \(\varphi_{m}(g)=\langle x_{m}|\pi(g)x_{m}\rangle\) converge to the constant function 1 in \(\mathcal{C}(\mathrm{G})\) endowed with the topology of compact convergence. (Consider a sequence \((\mathrm{C}_{m})_{m\in\mathbf{N}}\) of compact subsets whose union is equal to G, and take \(x_{m}\) satisfying the above condition for \((\mathrm{C}_{m},1/(m+1))\) .)
\item
  For every \(m \in N\) , there exists a positive measure \(\mu_{m}\) of total mass 1 on \(R^{n}\) whose Fourier transform is identified with the restriction of \(\varphi_{m}\) to H.
\item
  Suppose that \(\mu_{m}(\{0\}) = 0\) for all \(m \in N\) . We then identify \(\mu_{m}\) with a positive measure of total mass 1 on \(R^{n} - \{0\}\) , and we denote by \(\nu_{m}\) the image measure of \(\mu_{m}\) on projective space \(\mathbf{P}_{n-1}(\mathbf{R})\) . There exists a subsequence \((\nu_{n_{k}})_{k \in \mathbf{N}}\) converging weakly to a positive measure \(\nu\) of total mass 1 on \(\mathbf{P}_{n-1}(\mathbf{R})\) ; the measure \(\nu\) is \(\mathbf{SL}(n, \mathbf{R})\) -invariant. Deduce a contradiction.
\item
  Let \(m \in N\) be such that \(\mu_{m}(\{0\}) > 0\) . Let \(\psi\) be the restriction of \(\varphi_{m}\) to H. There exists c > 0 such that \(\psi = c + \widetilde{\psi}\) where \(\widetilde{\psi}\) is a function of positive type on \(R^{n}\) .
\item
  The representation \(\pi\) contains H-invariant vectors. (Use Exercise 2.) One says that the pair (G,H) has the relative ]T[ property.
\end{enumerate}

\begin{enumerate}
\def\labelenumi{\arabic{enumi})}
\setcounter{enumi}{17}
\tightlist
\item
  Let \(\mathrm{G} = \mathbf{SL}(2, \mathbf{R})\) . Let N and P denote the subgroups
\end{enumerate}

\[
\mathrm{N} = \Big \{\left( \begin{array}{c c} 1 & t \\ 0 & 1 \end{array} \right) | t \in \mathbf {R} \Big \}, \qquad \mathrm{P} = \Big \{\left( \begin{array}{c c} a & t \\ 0 & a ^ {- 1} \end{array} \right) | a \in \mathbf {R} _ {*}, t \in \mathbf {R} \Big \}
\]

of G. Let \(\pi\) be a unitary representation of G in a complex Hilbert space E. Suppose that there exists \(x \in \mathrm{E}\) such that \(\pi(n)x = x\) for all \(n \in \mathrm{N}\) .

\begin{enumerate}
\def\labelenumi{\alph{enumi})}
\tightlist
\item
  The homogeneous space \(\mathbf{SL}(2, \mathbf{R}) / \mathrm{N}\) is identified with \(\mathbf{R}^2 - \{0\}\) by the mapping \(g \mapsto ge_1\) where \(e_1\) is the first vector of the canonical basis of \(\mathbf{R}^2\) , and the homogeneous space \(\mathbf{SL}(2, \mathbf{R}) / \mathrm{P}\) is identified with the subspace \((\mathbf{R} - \{0\}) \times \{0\}\) of \(\mathbf{SL}(2, \mathbf{R}) / \mathrm{N}\) .
\item
  Every continuous function \(f\) on G such that \(f(ngn') = f(g)\) for all \((n, n') \in \mathrm{N}^2\) satisfies \(f(pgp') = f(g)\) for all \((p, p') \in \mathrm{P}^2\) . (Interpret \(f\) as a continuous function \(\varphi\) on \(\mathbf{R}^2 - \{0\}\) invariant under the natural action of N and verify that it is constant on \(\mathbf{SL}(2, \mathbf{R}) / \mathrm{P}\) .)
\end{enumerate}

\begin{enumerate}
\def\labelenumi{\alph{enumi})}
\setcounter{enumi}{2}
\tightlist
\item
  One has \(\pi(g)x = x\) for all \(g \in G\) .
\end{enumerate}

\begin{enumerate}
\def\labelenumi{\arabic{enumi})}
\setcounter{enumi}{18}
\tightlist
\item
  Let \(n \geqslant 2\) be an integer. We consider the real Lie group \(\mathrm{G} = \mathbf{SL}(n, \mathbf{R})\) . Let N and H denote the closed subgroups of G whose elements are of block form
\end{enumerate}

\[
\left( \begin{array}{c c} \operatorname{Id} _ {n - 1} & x \\ 0 & 1 \end{array} \right), \quad (x \in \mathbf {R} ^ {n - 1}), \qquad \left( \begin{array}{c c} g & 0 \\ 0 & 1 \end{array} \right), \quad (g \in \mathbf {S L} (n - 1, \mathbf {R}))
\]

respectively. The group N is identified with \(R^{n-1}\) and the group H is identified with \(\mathbf{SL}(n-1,\mathbf{R})\) .

\begin{enumerate}
\def\labelenumi{\alph{enumi})}
\item
  Let \(\pi\) be a unitary representation of G in a complex Hilbert space E. An element \(x\) of E is invariant under G if and only if it is invariant under N.
\item
  If \(n \geqslant 3\) , then the group G has property ]T[. (Combine the result of the previous question with the previous exercise.)
\item
  If \(n \geqslant 3\) , then the group \(\widetilde{G} = R^{n} \rtimes G\) has property ]T[. (Let \(\pi\) be a unitary representation such that \(1_{\widetilde{G}} \leqslant \pi\) ; there exists a vector \(x \neq 0\) invariant under the subgroup G of \(\widetilde{G}\) by the preceding question; show that the corresponding matrix coefficient is constant.)
\item
  If \(n \geqslant 3\) , the group \(\mathbf{SL}(n, \mathbf{Z})\) has property ]T[. (Use Exercise 7 of INT, VII, p.~116.)
\item
  Let \(n \geqslant 3\) . The group \(\mathbf{SL}(n, \mathbf{Z})\) is finitely generated. Let H be the family of normal subgroups of finite index of \(\mathbf{SL}(n, \mathbf{Z})\) ; it is infinite. For every finite generating subset S of \(\mathbf{SL}(n, \mathbf{Z})\) the family of Cayley graphs \((\mathcal{G}(\mathbf{SL}(n, \mathbf{Z})/\mathrm{H}, \mathrm{SH}/\mathrm{H}))_{\mathrm{H} \in \mathcal{H}}\) is a family of expander graphs (cf.~TA, 2, p.~221, exercise 7 and exercise 16 of IV, p.~323).
\end{enumerate}

\begin{enumerate}
\def\labelenumi{\arabic{enumi})}
\setcounter{enumi}{19}
\tightlist
\item
  Let G be a locally compact group and \(\mu\) a left Haar measure on G. We denote by \(\Delta\) the modulus of G. For any function f on G, we will denote by \(\widetilde{f} = \Delta f^{*}\) .
\end{enumerate}

A measure \(\nu\) on G is said to be of positive type if \(\langle \nu, \tilde{f} * f \rangle \geqslant 0\) for every function \(f \in \mathcal{K}(\mathrm{G})\) . A locally integrable function \(\varphi\) on G is said to be of positive type if \(\langle \varphi, f^* * f \rangle \geqslant 0\) for every function \(f \in \mathcal{K}(\mathrm{G})\) .

\begin{enumerate}
\def\labelenumi{\alph{enumi})}
\item
  Suppose that \(\nu\) is bounded. The measure \(\nu\) is of positive type if and only if \(\gamma(\nu) \geqslant 0\) .
\item
  Let \(\varphi\in\mathcal{K}(\mathrm{G})\) . The function \(\varphi\) is of positive type as a locally integrable function if and only if it is so according to Definition 4 of V, p.~442.
\item
  Let \(\nu\) be a measure of positive type on G. The formula
\end{enumerate}

\[
(f \mid g) _ {\nu} = \langle \nu , \widetilde {f} * g \rangle
\]

defines a positive Hermitian form on \(\mathcal{K}(\mathrm{G})\) . We denote by \(\mathrm{E}_{\nu}\) the separated completion Hilbert space of \(\mathcal{K}(\mathrm{G})\) equipped with this Hermitian form.

\begin{enumerate}
\def\labelenumi{\alph{enumi})}
\setcounter{enumi}{3}
\tightlist
\item
  For \(f \in \mathcal{K}(\mathrm{G})\) , the function \(\sigma(g)f \colon g \mapsto (\varepsilon_g * f)\Delta^{1/2}\) belongs to \(\mathcal{K}(\mathrm{G})\) and \(\| \sigma(g)f \|_{\nu} = \|f\|_{\nu}\) . The map \(g \mapsto \sigma(g)\) extends to a unitary representation of G in \(\mathrm{E}_{\nu}\) .
\end{enumerate}

\begin{enumerate}
\def\labelenumi{\alph{enumi})}
\setcounter{enumi}{4}
\tightlist
\item
  Let U be an open neighborhood of e in G, and let \(\nu\) be a positive-type measure on G such that there exists a real number \(c \geqslant 0\) so that
\end{enumerate}

\[
\langle \nu , f * \widetilde {f} \rangle \leqslant c \left(\int_ {\mathrm{G}} f d \mu\right) ^ {2}
\]

  if \(f \in \mathcal{K}(\mathrm{G})\) is positive and vanishes outside U. There exists a continuous function of positive type \(\varphi\) on G such that \(\nu = \varphi \Delta^{-1/2} \cdot \mu\) . (Let \(j\) be the canonical morphism from \(\mathcal{K}(\mathrm{G})\) into \(\mathrm{E}_{\nu}\) ; if \(\mathfrak{B}\) is a basis for the filter of neighborhoods of \(e\) whose elements are contained in U and \(f_{\mathrm{V}} \in \mathcal{K}_{+}(\mathrm{G})\) is a function with support in \(V \in B\) of integral 1, prove that \(j(f_{\mathrm{V}})\) converges weakly according to B to an element \(y_{0}\) of \(E_{\nu}\) such that

\[
(j (f) \mid y _ {0}) = \int_ {\mathrm{G}} f ^ {*} d \mu
\]

for \(f \in \mathcal{K}(\mathrm{G})\) ; setting \(\varphi(g) = (y_0 \mid \sigma(g)y_0)\) , verify that \(\varphi\) has the required properties.)

\begin{enumerate}
\def\labelenumi{\alph{enumi})}
\setcounter{enumi}{5}
\tightlist
\item
  Let \(\varphi_{1}\) be a continuous positive-definite function on G and \(\varphi_{2}\) a locally integrable function of positive type. If \(\varphi_{1}-\varphi_{2}\) is of positive type, then \(\varphi_{2}\) is equal locally almost everywhere to a continuous function of positive type. (Verify that one can apply the previous question.)
\end{enumerate}

\begin{enumerate}
\def\labelenumi{\arabic{enumi})}
\setcounter{enumi}{20}
\tightlist
\item
  We retain the notation of the preceding exercise.
\end{enumerate}

\begin{enumerate}
\def\labelenumi{\alph{enumi})}
\item
  Let \(f \in \mathcal{L}^{2}(G)\) that is of positive type as a locally integrable function. The partial operator on \(L^{2}(G)\) with domain \(\mathcal{K}(G)\) defined by \(\varphi \mapsto \gamma(\varphi)f\) is positive. We denote by \(\varrho(f)\) the Friedrichs extension of f (example 8 of IV, p.~295); it is a positive self-adjoint partial operator on \(L^{2}(G)\) .
\item
  Let \(f \in \mathcal{L}^2(\mathrm{G})\) . One says that \(f\) is sober if there exists a real number \(c \geqslant 0\) so that \(\| \boldsymbol{\gamma}(\varphi)f \| \leqslant c \| \varphi \|\) for all \(\varphi \in \mathcal{K}(\mathrm{G}) \subset \mathcal{L}^2(\mathrm{G})\) . We then denote by \(\varrho(f)\) the endomorphism of \(\mathrm{L}^2(\mathrm{G})\) which extends \(\varphi \mapsto \boldsymbol{\gamma}(\varphi)f\) . This definition coincides with the previous one if \(f\) is of positive type, and this condition then amounts to saying that \(\varrho(f)\) is positive.
\item
  Let \(f_{1}\) and \(f_{2}\) be sober functions in \(\mathcal{L}^{2}(G)\) . If \(\varrho(f_{1}) = \varrho(f_{2})\) , then \(f_{1}\) and \(f_{2}\) are equal in \(L^{2}(G)\) .
\end{enumerate}

\begin{enumerate}
\def\labelenumi{\alph{enumi})}
\setcounter{enumi}{3}
\tightlist
\item
  Let \(f \in \mathcal{L}^2(\mathrm{G})\) be of positive type. For all \(g \in \mathrm{G}\) , we have \(\varrho(f) \circ \boldsymbol{\gamma}(g) = \boldsymbol{\gamma}(g) \circ \varrho(f)\) . Moreover, for every \(h \in \operatorname{dom}(\varrho(f))\) , we have \(\varrho(f)h = h * f\) .
\end{enumerate}

\begin{enumerate}
\def\labelenumi{\alph{enumi})}
\setcounter{enumi}{4}
\tightlist
\item
  Let \(f_{1}\) and \(f_{2}\) be sober elements of \(\mathcal{L}^{2}(\mathrm{G})\) such that \(\varrho(f_{1})\) and \(\varrho(f_{2})\) commute. The function \(f_{1}*f_{2}\) is a continuous positive-definite function on G; it is sober, and we have
\end{enumerate}

\[
\varrho (f _ {1} * f _ {2}) = \varrho (f _ {1}) \varrho (f _ {2}), \quad \langle f _ {1} | f _ {2} \rangle \geqslant 0.
\]

\begin{enumerate}
\def\labelenumi{\alph{enumi})}
\setcounter{enumi}{5}
\tightlist
\item
  With the notations and hypotheses of the previous question, if \(f_{2} - f_{1}\) is of positive type, then
\end{enumerate}

\[
\left\| f _ {2} - f _ {1} \right\| ^ {2} \leqslant \left\| f _ {2} \right\| ^ {2} - \left\| f _ {1} \right\| ^ {2}.
\]

\begin{enumerate}
\def\labelenumi{\alph{enumi})}
\setcounter{enumi}{6}
\tightlist
\item
  Let \((f_n)_{n\in \mathbf{N}}\) be a sequence of sober positive-type elements of \(\mathcal{L}^2 (\mathrm{G})\) , pairwise commuting, and such that \(f_{n + 1} - f_n\) is of positive type for all \(n\in \mathbf{N}\) . If the sequence \((\| f\| _n)\) is bounded, then the sequence \((f_n)\) converges in \(\mathcal{L}^2 (\mathrm{G})\) .
\end{enumerate}

\begin{enumerate}
\def\labelenumi{\arabic{enumi})}
\setcounter{enumi}{21}
\tightlist
\item
  We retain the notation of the preceding exercise.
\end{enumerate}

Let \(f\) be a continuous function of positive type belonging to \(\mathcal{L}^2 (\mathrm{G})\) .

\begin{enumerate}
\def\labelenumi{\alph{enumi})}
\item
  There exists an increasing sequence \((p_{n})_{n\in\mathbf{N}}\) of polynomial functions on [0,1] which converges uniformly on [0,1] to the function \(t\mapsto\sqrt{t}\) .
\item
  Suppose that \(f\) is sober and that \(0 \leqslant f \leqslant 1\) in \(\mathcal{L}(\mathrm{L}^2(\mathrm{G}))\) . There exists a sequence \((h_n)\) in \(\mathcal{L}^2(\mathrm{G})\) satisfying the following conditions:
\end{enumerate}

\begin{enumerate}
\def\labelenumi{(\roman{enumi})}
\item
  Each \(h_n\) is sober and of positive type;
\item
  The function \(h_{n+1} - h_n\) is of positive type for all \(n \in \mathbf{N}\) ;
\item
  The endomorphisms \(\varrho(h_n)\) commute pairwise;
\item
  One has \(\varrho(h_n)^2 \leqslant \varrho(f)\) for all \(n\) .
\end{enumerate}

\begin{enumerate}
\def\labelenumi{\alph{enumi})}
\setcounter{enumi}{2}
\tightlist
\item
  The sequence \((h_n)_{n\in \mathbf{N}}\) converges in \(\mathcal{L}^2 (\mathrm{G})\) ; if \(h\) is a limit of \((h_n)\) , then \(h\) is of positive type and sober, and satisfies \(\varrho (h)^2 = \varrho (f)\) ; we have \(f = h*h\) .
\end{enumerate}

\begin{enumerate}
\def\labelenumi{\alph{enumi})}
\setcounter{enumi}{3}
\tightlist
\item
  We no longer assume that \(f\) is sober. For \(t \in \mathbf{R}\) , we denote by \(p_t\) the spectral projection of the partial self-adjoint operator \(\varrho(f)\) defined by \([- \infty, t]\) . We have \(p_t \circ \boldsymbol{\gamma}(g) = \boldsymbol{\gamma}(g) \circ p_t\) for all \(t \in \mathbf{R}\) and \(g \in G\) .
\end{enumerate}

\begin{enumerate}
\def\labelenumi{\alph{enumi})}
\setcounter{enumi}{4}
\tightlist
\item
  For \(t \in \mathbf{R}\) , set \(f_t = p_t(f)\) ; this is the class in \(\mathrm{L}^2(\mathrm{G})\) of a continuous positive-definite sober function in \(\mathcal{L}^2(\mathrm{G})\) (Check that \(\gamma(\varphi)f_t = \varrho(f)p_t(g)\) for \(g \in \mathcal{H}(\mathrm{G})\) , and deduce that \(\varrho(f_t) = \varrho(f) \circ p_t\) , then apply the results of Exercise 20.)
\end{enumerate}

\begin{enumerate}
\def\labelenumi{\alph{enumi})}
\setcounter{enumi}{5}
\tightlist
\item
  There exists \(h \in \mathcal{L}^2(\mathrm{G})\) such that \(f = h * h\) (Consider \(f_n = p_n(f)\) for \(n \in \mathbf{N}\) and apply part \(c\) ).)
\end{enumerate}

\begin{enumerate}
\def\labelenumi{\arabic{enumi})}
\setcounter{enumi}{22}
\tightlist
\item
  We denote \(G = R_{+}^{*}\) ; it is a locally compact commutative group and the measure \(\mu = dx/x\) is a Haar measure on G. The group G and the group iR are in duality via the mapping \((x, it) \mapsto x^{it}\) ; the dual measure of \(\mu\) is Lebesgue measure on iR (cf. Cor. 4 of II, p. 236). We denote by F the Fourier transform of G, and \(\overline{F}\) the Fourier cotransform of \(\widehat{G}\) .
\end{enumerate}

We denote by X the identity function on G.

\begin{enumerate}
\def\labelenumi{\alph{enumi})}
\item
  Let \(\nu\) be a positive measure on G. We denote by \(\mathrm{I}_{\nu}\) the set of \(\sigma \in \mathbf{R}\) such that the function \(\mathrm{X}^{\sigma}\) is \(\nu\) -integrable, and \(\mathrm{B}_{\nu}\) the set of \(s \in \mathbf{R}\) such that \(\mathcal{R}(s) \in \mathrm{I}_{\nu}\) . The set \(\mathrm{I}_{\nu}\) is an interval of \(\mathbf{R}\) and the function \(s \mapsto \nu(\mathrm{X}^s)\) is defined and holomorphic in the interior of \(\mathrm{B}_{\nu}\) .
\item
  Let \(f\) be a locally integrable function on G. We denote by \(\mathrm{I}_{f} = \mathrm{I}_{f\cdot \mu}\) and \(\mathrm{B}_f = \mathrm{B}_{f\cdot \mu}\) . For every \(s \in \mathrm{B}_f\) we set
\end{enumerate}

\[
\widehat {f} (s) = \int_ {\mathrm{G}} f (x) x ^ {s} d \mu (x).
\]

If there exists \(\delta > 0\) and \(c \geqslant 0\) such that \(|\widehat{f}(s)| \leqslant c(1 + |s|)^{-1 - \delta}\) for \(s \in \mathrm{B}_f\) , then for all \(\sigma\) in the interior of \(\mathrm{B}_f\) , we have

\[
f (x) = \frac {1}{2 \pi} \int_ {\mathbf {R}} \widehat {f} (\sigma + i t) x ^ {- \sigma - i t} d t.
\]

\begin{enumerate}
\def\labelenumi{\alph{enumi})}
\setcounter{enumi}{2}
\tightlist
\item
  Let \(f \in \mathcal{L}^{2}(\mathrm{G})\) . We assume that 0 belongs to the interior of \(I_{f}\) . We then have \(\overline{\mathcal{F}} f(it) = \widehat{f}(it)\) for almost every \(t \in \mathbf{R}\) .
\end{enumerate}

\begin{enumerate}
\def\labelenumi{\alph{enumi})}
\setcounter{enumi}{3}
\item
  Let \(f\) and \(g\) be measurable functions on \(\mathbf{G}\) . Let \(\sigma\) be a real number such that \(\mathrm{X}^{-2\sigma}f\in \mathrm{L}^2 (\mathrm{G})\) and \(\mathrm{X}^{2\sigma}g\in \mathrm{L}^2 (\mathrm{G})\) . We then have

\[
\int_ {\mathrm{G}} f g d \mu = \frac {1}{2 \pi} \int_ {\mathbf {R}} \widehat {f} (- \sigma - i t) \widehat {g} (\sigma + i t) d t.
\]

\end{enumerate}

\begin{enumerate}
\def\labelenumi{\alph{enumi})}
\setcounter{enumi}{4}
\tightlist
\item
  Let \(f\) and \(g\) be measurable functions on \(G\) . One assumes that
\end{enumerate}

\begin{enumerate}
\def\labelenumi{(\roman{enumi})}
\item
  There exists \(c \in \mathbf{R}\) such that \(|f(x)| \leqslant cx / (1 + x)\) for all \(x \in \mathbf{G}\) ;
\item
  One has ]0,1[ ⊂ I \(_{g}\) and for every σ such that 0 < σ < 1, the function X \(^{\sigma}\) g is bounded;
\item
  There exists \(c' \in \mathbf{R}\) such that \(|\widehat{g}(s)| \leqslant c'(1 + |s|)^{-2}\) for all \(s \in \mathbf{C}\) such that \(\mathcal{R}(s) \in ]0, 1[\) .
\end{enumerate}

Then we have

\[
\int_ {\mathrm{G}} f g d \mu = \frac {1}{2 \pi} \int_ {\mathbf {R}} \widehat {f} (- \sigma - i t) \widehat {g} (\sigma + i t) d t
\]

for all \(\sigma\in\mathbf{R}\) such that \(0<\sigma<1\) .

\begin{enumerate}
\def\labelenumi{\alph{enumi})}
\setcounter{enumi}{5}
\tightlist
\item
  Let \(\sigma \in \mathbf{R}\) and let \(f \in \mathcal{C}(\sigma + i\mathbf{R})\) . There exists a positive measure \(\nu\) on G such that \(\sigma \in \mathrm{I}_{\nu}\) and \(\widehat{\nu}(\sigma + it) = f(\sigma + it)\) for all \(t \in \mathbf{R}\) if and only if for every integer \(n \geqslant 1\) , every family \((t_i)_{0 \leqslant i \leqslant n}\) of complex numbers and any family \((z_i)_{0 \leqslant i \leqslant n}\) of complex numbers such that \(\mathcal{R}(z_i) = \sigma\) for all \(i\) , we have
\end{enumerate}

\[
\sum_ {i = 0} ^ {n} \sum_ {j = 0} ^ {n} \bar {t} _ {i} t _ {j} f \left(\frac {1}{2} (\overline {{z}} _ {i} + z _ {j})\right) \geqslant 0.
\]

\exercisesection{}

\begin{enumerate}
\def\labelenumi{\arabic{enumi})}
\tightlist
\item
  Let G be a compact topological group and let B be the family of characters of irreducible unitary representations of G.
\end{enumerate}

\begin{enumerate}
\def\labelenumi{\alph{enumi})}
\item
  The pair \((\mathrm{R}(\mathrm{G}),\mathcal{B})\) is a natural unital basal algebra (cf.~exercise 6 of III, p.~130).
\item
  Assume that G is finite. The algebra R(G) is then of finite rank; compute \(p_f(\pi)\) for every \(\pi \in \widehat{G}\), and identify the representation of G corresponding to the element r₀ defined in question j) of loc. cit.
\end{enumerate}

\begin{enumerate}
\def\labelenumi{\arabic{enumi})}
\setcounter{enumi}{1}
\tightlist
\item
  Let \(n \geqslant 1\) and \(p \in [1, +\infty]\) . We equip \(\mathbf{R}^n\) with the norm
\end{enumerate}

\[
\| x \| = \left(\sum_ {i = 1} ^ {n} | x _ {i} | ^ {p}\right) ^ {1 / p} \text {   if   } p \neq + \infty ,
\]

and \(\|x\| = \sup_{i}|x_{i}|\) if \(p = +\infty\) . Let G be the group of isometries of the metric space \(R^{n}\) equipped with this norm.

\begin{enumerate}
\def\labelenumi{\alph{enumi})}
\item
  Let \(\sigma\) be a permutation of \(\{1,\ldots,n\}\) and \(\varepsilon\in\{-1,1\}^{n}\) . Let \((e_{i})\) be the canonical basis of \(\mathbf{R}^{n}\) . The unique endomorphism of \(\mathbf{R}^{n}\) such that \(e_{i}\mapsto\varepsilon_{i}e_{\sigma(i)}\) belongs to G. The set of isometries of this form is a subgroup W of G; it is an extension of \(\mathfrak{S}_{n}\) by \((\mathbf{Z}/2\mathbf{Z})^{n}\) .
\item
  The group G is compact.
\item
  The action of the group G on \(R^{n}\) is irreducible; the group G is contained in the orthogonal group of the quadratic form \(\sum x_{i}^{2}\) .
\item
  If \(p \neq 2\) , then G = W.
\end{enumerate}

\begin{enumerate}
\def\labelenumi{\arabic{enumi})}
\setcounter{enumi}{2}
\tightlist
\item
  We recall that a group \(\Gamma\) is said to be residually finite if the intersection of the finite-index subgroups of \(\Gamma\) is reduced to \(\{e\}\) (A, I, p.~129, exerc. 5, b)).
\end{enumerate}

For a group \(\Gamma\) , we denote by \(\widehat{\Gamma}\) the projective limit \(\widehat{\Gamma} = \varprojlim \Gamma / H\) , where \(H\) ranges over the set of finite-index normal subgroups of \(\Gamma\) ordered by inclusion. We say that \(\widehat{\Gamma}\) is the profinite completion of \(\Gamma\) .

\begin{enumerate}
\def\labelenumi{\alph{enumi})}
\item
  A group \(\Gamma\) is residually finite if and only if, for every \(g \in \Gamma\) different from the identity element, there exists a finite group F and a homomorphism \(\varphi: \Gamma \to F\) such that \(\varphi(g) \neq e\) .
\item
  Show that \(\Gamma\) is residually finite if and only if the canonical homomorphism from \(\Gamma\) into \(\widehat{\Gamma}\) is injective.
\item
  Let K be a field. Every finitely generated subgroup (cf.~A, I, p.~145, exerc. 16) of \(\mathbf{GL}(n,\mathrm{K})\) is residually finite ("Mal'cev's theorem"). (Show that there exists an integer \(m\geqslant 0\) , a ring A equal to \(\mathbf{Z}\) or to a finite field and an injective homomorphism of such a subgroup into \(\mathbf{GL}(m,\mathrm{R})\) , where R is of the form \(\mathrm{A}[\mathrm{X}_1,\dots ,\mathrm{X}_r,1 / f]\) .)
\item
  Let \(\Gamma\) be a finitely generated discrete group. Suppose that there exists an injective homomorphism from \(\Gamma\) into a compact topological group. For every \(\gamma \in \Gamma\) different from the identity element, there exists a finite-dimensional unitary irreducible representation of \(\Gamma\) such that the image of \(\gamma\) by \(\pi\) is not the identity map.
\item
  A finitely generated group is residually finite if and only if it is isomorphic to a subgroup of a compact group.
\end{enumerate}

¶ 4) Let G be a group and let A and B be subgroups of G. Let \(\varphi\colon A\to B\) be an isomorphism.

Let \(r'\) be a family in the free group F(G) such that G is isomorphic to the group F(G, \(r'\) ) (A, I, p.~87, n°6). One calls the HNN extension of G according to \(\varphi\) the group \(G^{\varphi} = F(X, r)\) where X is the sum of the sets G, A, B and an element t, and where r is the sum family of \(r'\) and of the relators \(t^{-1}at = \varphi(a)\) for all \(a \in A\) .

We denote by \(j\colon G\to G^{\varphi}\) the unique homomorphism such that \(j(g)\) is the element represented by \(g\) in \(G^{\varphi}\) .

\begin{enumerate}
\def\labelenumi{\alph{enumi})}
\tightlist
\item
  Let \(k \geqslant 1\) an integer and let \((g_{0}, \varepsilon_{1}, g_{1}, \ldots, \varepsilon_{k}, g_{k})\) be a finite family such that \(g_{i} \in G\) and \(\varepsilon_{i} \in \{-1, 1\}\) for every i. Suppose that there does not exist an integer j such that \(1 \leqslant j \leqslant k - 1\) and
\end{enumerate}

\[
(\varepsilon_ {j}, \varepsilon_ {j + 1}) = (- 1, 1), \quad g _ {i} \in A
\]

or

\[
(\varepsilon_ {j}, \varepsilon_ {j + 1}) = (1, - 1), \quad g _ {i} \in \mathrm{B}.
\]

Then the element

\[
j (g _ {0}) t ^ {\varepsilon_ {1}} j (g _ {1}) \dots t ^ {\varepsilon_ {k}} j (g _ {k})
\]

of G \(^{φ}\) is not the identity element ("Britton's Lemma".)

\begin{enumerate}
\def\labelenumi{\alph{enumi})}
\setcounter{enumi}{1}
\item
  The homomorphism j is injective.
\item
  Let \(a\) and \(b\) be distinct symbols. Let G be the group \(\mathrm{F}(\{a,b\}, a^{-1}b^{2}ab^{-3})\) . We denote by \(b_{1} = a^{-1}ba\) . Let H be a normal subgroup of finite index of G and \(\pi: \mathrm{G} \to \mathrm{G}/\mathrm{H}\) the canonical projection. One has \([b_{1}, b] \in \mathrm{Ker}(\pi)\) .
\item
  The group G admits no faithful finite-dimensional representation. (Show that G is not residually finite.)
\end{enumerate}

\begin{enumerate}
\def\labelenumi{\arabic{enumi})}
\setcounter{enumi}{4}
\tightlist
\item
  Let G be a compact group and \(\varrho\) a faithful finite-dimensional unitary representation of G. Let \(\pi \in \widehat{G}\) .
\end{enumerate}

\begin{enumerate}
\def\labelenumi{\alph{enumi})}
\tightlist
\item
  There exist integers a and b such that \(\pi\) is isomorphic to a subrepresentation of \(\varrho^{\otimes a} \otimes \overline{\varrho}^{\otimes b}\) . (One may use, for example, the Weierstrass--Stone theorem.)
\end{enumerate}

We denote by \(w(\pi)\) the minimum of \(a + b\) when \((a, b)\) ranges over the set of integers satisfying this property.

\begin{enumerate}
\def\labelenumi{\alph{enumi})}
\setcounter{enumi}{1}
\tightlist
\item
  Let \(\Lambda\) be the weight lattice of \(G^{0}\) and let q be a norm on \(\Lambda\otimes \mathbf{R}\) . There exist real numbers C and D such that for every representation \(\pi\in\widehat{G}\) whose restriction to \(G^{0}\) has highest weights \((\lambda_{1},\ldots,\lambda_{k})\) , we have
\end{enumerate}

\[
w (\pi) \leqslant \mathrm{C} \sup (q (\lambda_ {1}), \dots , q (\lambda_ {k})) + \mathrm{D}.
\]

(Consider first the case of a connected semisimple group and apply a) to a suitably chosen finite family of representations of G).

\begin{enumerate}
\def\labelenumi{\arabic{enumi})}
\setcounter{enumi}{5}
\tightlist
\item
  Let G be a compact group and let \(\varrho\) be an irreducible representation of G. There exists a central function \(f \in \Theta(G^{\sharp})\) such that
\end{enumerate}

\[
f = \sum_ {\pi \in \widehat {\mathrm{G}}} c (\pi) \chi_ {\pi}
\]

where the integers \(c(\pi)\) satisfy

\begin{enumerate}
\def\labelenumi{(\roman{enumi})}
\item
  \(c(\pi) \geqslant 0\) for all \(\pi \in \widehat{G}\) ;
\item
  \(c(1) \leqslant c(\varrho)\) , and the inequality is strict except possibly if \(\varrho\) is of dimension 1 and satisfies \(\varrho^{2} = 1\) ;
\item
  \(\mathcal{R}(f) \geqslant 0.\)
\end{enumerate}

¶ 7) For every integer \(n \geqslant 1\) , we denote by \(\mu_n\) the positive measure on \(\mathbf{C}\) image of the normalized Haar measure of the group \(\mathbf{U}(\mathbf{C}^n) \subset \mathrm{GL}(\mathbf{C}^n)\) under the trace map \(g \mapsto \mathrm{Tr}(g)\) ; it is a compactly supported measure of total mass 1.

\begin{enumerate}
\def\labelenumi{\alph{enumi})}
\tightlist
\item
  Let \(a\) and \(b\) be natural integers such that \(a + b \leqslant n\) . We have
\end{enumerate}

\[
\int_ {\mathbf {C}} z ^ {a} \overline {{z}} ^ {b} d \mu_ {n} (z) = \left\{ \begin{array}{l l} 0 & \text{if } a \neq b \\ a! & \text{if } a = b. \end{array} \right.
\]

\begin{enumerate}
\def\labelenumi{\alph{enumi})}
\setcounter{enumi}{1}
\tightlist
\item
  The sequence of measures \((\mu_{n})_{n\geqslant1}\) converges narrowly to the Gaussian measure on the real vector space C whose covariance is \(\mathrm{Q}(z)=|z|^{2}\) .
\end{enumerate}

\begin{enumerate}
\def\labelenumi{\arabic{enumi})}
\setcounter{enumi}{7}
\item
  Let G be a compact group and let \(\varrho_{1}\) and \(\varrho_{2}\) be irreducible unitary representations of G. Let \(n \geqslant 2\) be an integer. If \(S^{n}\varrho_{1}\) is isomorphic to \(S^{n}\varrho_{2}\) , then there exists a subgroup H of finite index of G such that the restriction of \(\varrho_{1}\) to H is isomorphic to the restriction of \(\varrho_{2}\) to H.
\item
  Let \(k\) be a finite field with 5 elements and \(\mathrm{G} = \mathbf{S}\mathbf{L}(2,k)\) .
\end{enumerate}

\begin{enumerate}
\def\labelenumi{\alph{enumi})}
\tightlist
\item
  The conjugacy classes of the group G are represented by e, -e and the seven elements
\end{enumerate}

\[
g _ {3} = \left( \begin{array}{c c} 1 & 1 \\ 0 & 1 \end{array} \right), \qquad g _ {4} = \left( \begin{array}{c c} 1 & 2 \\ 0 & 1 \end{array} \right), \qquad g _ {5} = \left( \begin{array}{c c} - 1 & 1 \\ 0 & - 1 \end{array} \right), \qquad g _ {6} = \left( \begin{array}{c c} - 1 & 2 \\ 0 & - 1 \end{array} \right)
\]

\[
g _ {7} = \left( \begin{array}{c c} 2 & 0 \\ 0 & - 2 \end{array} \right), \qquad g _ {8} = \left( \begin{array}{c c} 0 & 1 \\ - 1 & 1 \end{array} \right), \qquad g _ {9} = \left( \begin{array}{c c} 0 & 1 \\ - 1 & - 1 \end{array} \right).
\]

Determine the cardinality of each conjugacy class.

\begin{enumerate}
\def\labelenumi{\alph{enumi})}
\setcounter{enumi}{1}
\tightlist
\item
  Let \(\alpha = \frac{1}{2}(1 + \sqrt{5})\) . There exists a faithful irreducible unitary representation \(\pi\) of dimension 2 of G whose character \(\chi\) satisfies \(\chi(-e) = -2\) and
\end{enumerate}

\[
\begin{array}{c} \chi (g _ {4}) = \alpha , \quad \chi (g _ {6}) = - \alpha , \quad \chi (g _ {3}) = - \alpha - 1, \quad \chi (g _ {5}) = \alpha + 1 \\ \chi (g _ {7}) = 0, \quad \chi (g _ {8}) = 1, \quad \chi (g _ {9}) = - 1. \end{array}
\]

(One may consider the restriction to G of the irreducible representations of dimension 4 of \(\mathbf{GL}(2,k)\), cf.~A, VIII, p.~422, exercise 27.)

\begin{enumerate}
\def\labelenumi{\alph{enumi})}
\setcounter{enumi}{2}
\tightlist
\item
  One has \(\mathrm{M}_4(\pi) = 2\) .
\end{enumerate}

\begin{enumerate}
\def\labelenumi{\arabic{enumi})}
\setcounter{enumi}{9}
\tightlist
\item
  \begin{enumerate}
  \def\labelenumii{\alph{enumii})}
  \tightlist
  \item
    There exists an irreducible representation \(\pi\) of orthogonal type of the alternating group \(\mathfrak{A}_5\) such that \(\mathrm{M}_4(\pi) = 3\) .
  \end{enumerate}
\end{enumerate}

\begin{enumerate}
\def\labelenumi{\alph{enumi})}
\setcounter{enumi}{1}
\tightlist
\item
  Let R be a root system of type \(E_{6}\) (resp. of type \(E_{7}\), of type \(E_{8}\) ) and W the Weyl group of R. Let \(\pi\) be the geometric representation of W (LIE, V, § 4, no. 3). It is irreducible, of orthogonal type, and satisfies \(M_{4}(\pi)=3\) .
\end{enumerate}

(One may use software such as GAP or MAGMA to determine the list of conjugacy classes of W and their images under \(\pi\) .)

\begin{enumerate}
\def\labelenumi{\arabic{enumi})}
\setcounter{enumi}{10}
\item
  Let G be a compact group and \(\pi\) a finite-dimensional unitary representation of G in a Hilbert space E. If \(\mathrm{M}_{4}(\pi) \leqslant 5\) then the representation \(\pi\) is irreducible.
\item
  Let G be a connected compact semisimple Lie group whose complexified Lie algebra is isomorphic to the simple Lie algebra \(e_{6}\) .
\end{enumerate}

\begin{enumerate}
\def\labelenumi{\alph{enumi})}
\item
  There exists a faithful irreducible representation \(\varrho\) of G of dimension 27; such a representation is of complex type.
\item
  The fourth moment of the representation \(\varrho\) is equal to 3.
\end{enumerate}

\begin{enumerate}
\def\labelenumi{\arabic{enumi})}
\setcounter{enumi}{12}
\tightlist
\item
  Let G be a connected compact semisimple Lie group and let g be its Lie algebra. Suppose there exists a faithful irreducible representation \(\varrho\) of dimension 27 of G. We assume that the fourth moment of \(\varrho\) is equal to 3 and that the representation \(\varrho\) is of complex type.
\end{enumerate}

\begin{enumerate}
\def\labelenumi{\alph{enumi})}
\item
  The Lie algebra g is simple.
\item
  The Lie algebra g is isomorphic to \(\mathfrak{e}_{6}\) .
\end{enumerate}

\begin{enumerate}
\def\labelenumi{\arabic{enumi})}
\setcounter{enumi}{13}
\item
  Let G be a connected compact semisimple Lie group. Suppose there exists a faithful irreducible representation \(\varrho\) of dimension 27 of G. The complexified Lie algebra of G is isomorphic to \(\mathfrak{e}_6\) if and only if there exists an irreducible subrepresentation of dimension 78 in the representation \(\operatorname{End}(\varrho)\) .
\item
  Let \(q\) be a power of an odd prime number and F a finite field with \(q\) elements. Consider the Heisenberg group H \(_3\) (F) of matrices of size 3 with coefficients in F that are upper triangular and unipotent. For \((a, b, c) \in \mathrm{F}^3\) , we denote
\end{enumerate}

\[
\sigma (a, b, c) = \left( \begin{array}{c c c} 1 & a & c \\ 0 & 1 & b \\ 0 & 0 & 1 \end{array} \right) \in \mathrm{H} _ {3} (\mathrm{F}).
\]

The map \(\sigma\) is therefore a bijection.

\begin{enumerate}
\def\labelenumi{\alph{enumi})}
\item
  The center C of \(\mathrm{H}_{3}(\mathrm{F})\) is identified with F by the group homomorphism sending \(\sigma(a,b,c)\) to c.
\item
  There exist \(q^{2}\) one-dimensional unitary characters of \(\mathrm{H}_{3}(\mathrm{F})\) ; these characters are trivial on C.
\item
  Let X be the subgroup of index q of F consisting of the elements \(\sigma(0,b,c)\) for \((b,c)\in\mathrm{F}^{2}\) . The subgroup X is identified with \(F^{2}\) . Let \(\psi_{1}\) and \(\psi_{2}\) be unitary characters of F, and set \(\psi(\sigma(0,b,c))=\psi_{1}(b)\psi_{2}(c)\) . The unitary representation \(\varrho\) of \(\mathrm{H}_{3}(\mathrm{F})\) induced by the character \(\psi\) of X is irreducible if and only if \(\psi_{2}\) is nontrivial; it depends only on \(\psi_{2}\) up to isomorphism, and its central character is equal to \(\psi_{2}\) . The support of the character of \(\varrho\) is equal to C.
\item
  The irreducible representations constructed in the two preceding questions are all the irreducible unitary representations of \(\mathrm{H}_{3}(\mathrm{F})\) .
\end{enumerate}

\begin{enumerate}
\def\labelenumi{\alph{enumi})}
\setcounter{enumi}{4}
\tightlist
\item
  Let \(\varrho\) be one of the irreducible unitary representations of dimension q of \(\mathrm{H}_{3}(\mathrm{F})\) . The representation \(\varrho\) is faithful and satisfies \(\mathrm{M}_{4}(\varrho)=q^{2}\) .
\end{enumerate}

\begin{enumerate}
\def\labelenumi{\arabic{enumi})}
\setcounter{enumi}{15}
\tightlist
\item
  Let \(r \geqslant 1\) be an integer and let G be a connected compact subgroup of \(\mathbf{U}(\mathbf{C}^{r})\) acting irreducibly on \(C^{r}\) . We denote by D the subgroup of diagonal matrices in \(\mathbf{U}(\mathbf{C}^{r})\) .
\end{enumerate}

\begin{enumerate}
\def\labelenumi{\alph{enumi})}
\tightlist
\item
  Let A be a subgroup of D contained in the normalizer of G in \(\mathbf{U}(\mathbf{C}^{r})\) . Let S be the subset of \(\widehat{A}\) formed by the diagonal coefficients \(a = (a_{i,j}) \mapsto a_{i,i}\) , where \(1 \leqslant i \leqslant r\) . We assume that S contains r elements and that the equation
\end{enumerate}

\[
\chi_ {1} \chi_ {2} ^ {- 1} = \chi_ {3} \chi_ {4} ^ {- 1}
\]

with \(\chi_{i} \in S\) admits only the solutions of the form \((\chi_{1}, \chi_{1}, \chi_{3}, \chi_{3})\) and \((\chi_{1}, \chi_{2}, \chi_{1}, \chi_{2})\) (one then says that \(S\) is a Sidon set).

If G is semisimple, prove that \(\mathbf{G} = \mathbf{S}\mathbf{U}(\mathbf{C}^{r})\) (Consider the representation of A on \(\operatorname{End}(\mathbf{C}^{r})\) , and deduce from the hypothesis that the complexified Lie algebra of G is stable under conjugation by D.)

\begin{enumerate}
\def\labelenumi{\alph{enumi})}
\setcounter{enumi}{1}
\tightlist
\item
  The group G contains \(\mathbf{SU}(\mathbf{C}^{r})\) if and only if there exists an element \(g \in G\) such that the eigenvalues of g are distinct and form a Sidon set in \(C^{*}\) .
\end{enumerate}

\begin{enumerate}
\def\labelenumi{\arabic{enumi})}
\setcounter{enumi}{16}
\tightlist
\item
  Let G be a compact group and H a closed subgroup of G. Let \(\varrho\) be a unitary representation of H in a finite-dimensional Hilbert space E.
\end{enumerate}

\begin{enumerate}
\def\labelenumi{\alph{enumi})}
\tightlist
\item
  Let \(\pi \in \widehat{\mathbf{G}}\) . The space \(\mathrm{Hom}_{\mathrm{G}}(\pi, \mathrm{Ind}_{\mathrm{H}}^{\mathrm{G}}(\varrho))\) is of finite dimension.

\end{enumerate}

\begin{enumerate}
\def\labelenumi{\alph{enumi})}
\setcounter{enumi}{1}
\item
  For every \(u \in \mathrm{Hom}_{\mathrm{G}}(\pi, \mathrm{Ind}_{\mathrm{H}}^{\mathrm{G}}(\varrho))\) , the image of \(u\) is contained in the image of \(\mathcal{C}(\mathrm{G};\mathrm{E})\) in \(\mathrm{L}^2 (\mathrm{G};\mathrm{E})\) .
\item
  There exists a linear map \(\Phi\) from the space \(\mathrm{Hom}_{\mathrm{G}}(\pi, \mathrm{Ind}_{\mathrm{H}}^{\mathrm{G}}(\varrho))\) into the space \(\mathrm{Hom}_{\mathrm{H}}(\mathrm{Res}_{\mathrm{H}}^{\mathrm{G}}(\pi), \varrho)\) such that, for every \(u \in \mathrm{Hom}_{\mathrm{G}}(\pi, \mathrm{Ind}_{\mathrm{H}}^{\mathrm{G}}(\varrho))\) and every \(x \in \mathrm{E}_{\pi}\) , \(\Phi(u)(x)\) is the image at \(e\) of the continuous function representing \(u(x)\) .
\item
  The linear map \(\Phi\) is an isomorphism.
\item
  Suppose that \(\varrho\) is irreducible. The multiplicity of \(\pi\) in \(\mathrm{Ind}_{\mathrm{H}}^{\mathrm{G}}(\varrho)\) is equal to the multiplicity of \(\varrho\) in \(\mathrm{Res}_{\mathrm{H}}^{\mathrm{G}}(\pi)\) (Frobenius reciprocity formula, cf. A, VIII, p. 392).
\end{enumerate}

\begin{enumerate}
\def\labelenumi{\arabic{enumi})}
\setcounter{enumi}{17}
\item
  Assume that G is compact. Let \(\varrho\) and \(\varrho'\) be unitary representations of G. The representation \(\varrho\) is weakly contained in \(\varrho'\) (cf. exercise 11 of V, p. 497) if and only if for every irreducible unitary representation \(\pi\) of G, the condition \(m_{\pi}(\varrho) \geqslant 1\) implies that \(m_{\pi}(\varrho') \geqslant 1\) . (Reduce to the case where \(\varrho'\) is irreducible; for every normalized diagonal matrix coefficient \(f\) of \(\varrho\) , prove by applying Exercise 13 of V, p.~498 that there exists a diagonal matrix coefficient of \(\varrho'\) which is not orthogonal to \(f\) , and conclude using the orthogonality relations.)
\item
  Let G be a compact group equipped with its normalized Haar measure \(\mu\) and let \(f \in \mathcal{C}(\mathrm{G})\) a continuous function on G. There exists an irreducible representation \(\pi \in \widehat{G}\) such that \(f = \dim(\pi)^{-1}\chi_{\pi}\) if and only if
\end{enumerate}

\[
\int_ {\mathrm{G}} f (x g x ^ {- 1} h) d \mu (x) = f (g) f (h)
\]

for all \((g,h)\in \mathrm{G}^2\)

\begin{enumerate}
\def\labelenumi{\arabic{enumi})}
\setcounter{enumi}{19}
\tightlist
\item
  Let H be a locally compact topological group equipped with a left Haar measure \(\nu\) . Let G be a compact subgroup of H. Let \(\varrho\) be an irreducible unitary representation of H in a Hilbert space E.
\end{enumerate}

We assume that the multiplicity of \(\pi\) in the restriction of \(\varrho\) to G is finite for every irreducible representation \(\pi \in \widehat{G}\) . We identify a function on G with a compactly supported measure on H.

\begin{enumerate}
\def\labelenumi{\alph{enumi})}
\item
  For every \(\pi \in \widehat{\mathrm{G}}\) , the endomorphism \(\varrho(f * \chi_{\pi})\) of E is of finite rank.
\item
  For every \(f \in \mathrm{L}^{1}(\mathrm{H})\) , the endomorphism \(\varrho(f)\) is compact.
\end{enumerate}

\begin{enumerate}
\def\labelenumi{\arabic{enumi})}
\setcounter{enumi}{20}
\tightlist
\item
  Let \(n \geqslant 1\) be an integer.
\end{enumerate}

\begin{enumerate}
\def\labelenumi{\alph{enumi})}
\tightlist
\item
  For every integer \(r > n^{2}\) and every family \((a_{1}, \ldots, a_{r})\) of elements of \(\mathcal{L}(\mathbf{C}^{n})\) , we have
\end{enumerate}

\[
\sum_ {\sigma \in \mathfrak {S} _ {r}} \varepsilon (\sigma) a _ {\sigma (1)} \dots a _ {\sigma (r)} = 0.\tag{4}
\]

\begin{enumerate}
\def\labelenumi{\alph{enumi})}
\setcounter{enumi}{1}
\tightlist
\item
  We denote by \(r(n)\) the smallest integer \(r \geqslant 1\) such that formula (4) is valid for every family \((a_i)_{1 \leqslant i \leqslant r}\) . We have \(r(n) \geqslant r(n-1) + 2\) (Consider a family \((a_1, \ldots, a_{r(n-1)-1})\) in \(\mathcal{L}(\mathbf{C}^{n-1})\) such that the expression above, denoted b, is nonzero; interpreting the elements of \(\mathcal{L}(\mathbf{C}^n)\) as matrices, choose i and j such that \(b_{i,j}\) is nonzero; view \(a_i\) and b as elements of \(\mathcal{L}(\mathbf{C}^n)\) by adding a row and a column of zeros, then define matrices \(a_{r(n-1)}\) and \(a_{r(n-1)+1}\) so that expression (4) for \((a_1, \ldots, a_{r(n-1)+1})\) is nonzero.)
\end{enumerate}

\begin{enumerate}
\def\labelenumi{\arabic{enumi})}
\setcounter{enumi}{21}
\tightlist
\item
  Let \(r \geqslant 1\) be an integer and let \(\mathrm{H} \subset \mathbf{GL}(r, \mathbf{C})\) be a closed subgroup.
\end{enumerate}

\begin{enumerate}
\def\labelenumi{\alph{enumi})}
\item
  The set of functions on H of the form \(g \mapsto \langle \varrho(g)x, \lambda \rangle\) , where \(\varrho\) is a continuous linear representation of H in a finite-dimensional complex vector space E, \(x \in E\) and \(\lambda \in E'\) , is a dense involutive subalgebra of \(\mathcal{C}(\mathrm{H})\) .
\item
  For any function \(f \in \mathcal{K}(\mathrm{H})\) nonzero, there exists a continuous linear representation \(\varrho\) of H in a finite-dimensional complex vector space such that \(\varrho(f)\) is nonzero.
\item
  Assume that H is a connected semisimple real Lie group. For every \(f \in \mathcal{K}(H)\) nonzero, there exists an irreducible continuous linear representation \(\varrho\) of H in a finite-dimensional complex vector space such that \(\varrho(f)\) is nonzero.
\end{enumerate}

\begin{enumerate}
\def\labelenumi{\arabic{enumi})}
\setcounter{enumi}{22}
\tightlist
\item
  We retain the hypotheses and notation of the preceding exercise, and assume that H is a connected semisimple real Lie group.
\end{enumerate}

Let G be a compact subgroup of H. Assume that there exists a connected solvable closed subgroup N of H such that H = GN. (For every semisimple real Lie group, this condition is satisfied when G is a maximal compact subgroup of H; see for example INT, VII, p.~91, § 3, n° 3, prop. 7 when G = SL(n, R).)

Let \(\pi \in \widehat{\mathbf{G}}\) and \(e_{\pi} = \dim (\pi)\overline{\chi}_{\pi}\) . Let \(\varrho\) be the irreducible representation of H given by the last question of the preceding exercise; we denote by E the space of \(\varrho\) .

\begin{enumerate}
\def\labelenumi{\alph{enumi})}
\item
  There exists a cyclic vector for the restriction of \(\varrho\) to G. (Use Lie's theorem, LIE, III, p.~241, cor., to find a nonzero \(x \in E\) such that C \(x\) is invariant under the action of N.) The multiplicity of \(\pi\) in the restriction of \(\varrho\) to G is at most equal to \(\dim(\pi)\) .
\item
  Let \(r\) be the integer \(r(\dim (\pi)^2)\) defined in Exercise 22; the identity
\end{enumerate}

\[
\sum_ {\sigma \in \mathfrak {S} _ {r}} \varepsilon (\sigma) a _ {\sigma (1)} \dots a _ {\sigma (r)} = 0
\]

is valid for every family \(a_{i}=e_{\pi}*g_{i}*e_{\pi}\) of elements of the algebra \(\mathcal{M}_{c}(\mathrm{H})\) of measures with compact support on H, where \(g_{i}\in\mathcal{K}(\mathrm{H})\) for \(1\leqslant i\leqslant r\) .

\begin{enumerate}
\def\labelenumi{\alph{enumi})}
\setcounter{enumi}{2}
\tightlist
\item
  Let \(\varpi\) be an irreducible unitary representation of H in a Hilbert space F. Let \(F_{\pi}\) be the \(\pi\)-isotypic component of the restriction of \(\varpi\) to G and let \(p\) be the orthogonal projector of F with image \(F_{\pi}\) . For every family \((u_i)_{1 \leqslant i \leqslant r}\) in \(\mathcal{L}(F_{\pi})\) , we have
\end{enumerate}

\[
\sum_ {\sigma \in \mathfrak {S} _ {r}} \varepsilon (\sigma) u _ {\sigma (1)} \dots u _ {\sigma (r)} = 0
\]

  (use Exercise 1 of V, p.~486).

\begin{enumerate}
\def\labelenumi{\alph{enumi})}
\setcounter{enumi}{3}
\tightlist
\item
  The multiplicity of \(\pi\) in the restriction of \(\varpi\) to G is at most equal to \(\dim(\pi)\) . (Use the result of Exercise 22 of V, p.~510.)
\end{enumerate}

\begin{enumerate}
\def\labelenumi{\arabic{enumi})}
\setcounter{enumi}{23}
\tightlist
\item
  A topological group G is said to be minimally almost periodic if every continuous homomorphism from G into a compact group is trivial.
\end{enumerate}

\begin{enumerate}
\def\labelenumi{\alph{enumi})}
\item
  A topological group G is minimally almost periodic if and only if G admits no non-trivial finite-dimensional unitary representation. If G is commutative, then G is minimally almost periodic if and only if every continuous unitary character of G is trivial.
\item
  Let G be a topological group and \(g \in G\) such that, for every integer \(k \geqslant 1\) , there exists an integer \(m \geqslant 1\) such that \(g^{km}\) is conjugate to g. Then one has \(\varrho(g) = 1_{\mathrm{E}}\) for every unitary representation \(\varrho\) of finite dimension of G in a Hilbert space E.
\item
  Let k be a subfield of C. The group \(\mathbf{SL}(2,k)\) endowed with the discrete topology is minimally almost periodic. (Verify that the preceding question applies to \(\begin{pmatrix}1 & z \\ 0 & 1\end{pmatrix}\) for all \(z \in k\) .)
\end{enumerate}

\begin{enumerate}
\def\labelenumi{\arabic{enumi})}
\setcounter{enumi}{24}
\tightlist
\item
  Let G be a locally compact group endowed with a left Haar measure \(\nu\) and \(K \subset G\) a compact subgroup of G endowed with the normalized Haar measure \(\mu\) .
\end{enumerate}

We denote by \(\mathcal{H}(\mathrm{G},\mathrm{K})\) the space of functions \(\varphi \in \mathcal{H}(\mathrm{G})\) such that \(\varphi (k_1gk_2) = \varphi (g)\) for all \(g\in \mathrm{G}\) and every \((k_{1},k_{2})\in \mathrm{K}\times \mathrm{K}\) .

The space \(\mathcal{H}(\mathrm{G},\mathrm{K})\) is a unital subalgebra of the algebra \(\mathcal{M}_c(\mathrm{G})\) ; one says that \((\mathrm{G},\mathrm{K})\) is a Gelfand pair if the algebra \(\mathcal{H}(\mathrm{G},\mathrm{K})\) is commutative. \(a)\) The map \(u^{\mathrm{K}}\) which to \(f\) associates the function \(f^{\mathrm{K}}\) defined by

\[
f ^ {\mathrm{K}} (g) = \int_ {\mathrm{K} \times \mathrm{K}} f (k _ {1} g k _ {2}) d (\mu \otimes \mu) (k _ {1}, k _ {2})
\]

is a projector in \(\mathcal{K}(\mathrm{G})\) whose image is \(\mathcal{H}(\mathrm{G},\mathrm{K})\) .

\begin{enumerate}
\def\labelenumi{\alph{enumi})}
\setcounter{enumi}{1}
\tightlist
\item
  Let \(f\in \mathcal{K}(\mathrm{G})\) . We have
\end{enumerate}

\[
\int_ {\mathrm{G}} f (g) d \nu (g) = \int_ {\mathrm{G}} f ^ {\mathrm{K}} (g) d \nu (g)
\]

and \((f\varphi)^{\mathrm{K}} = f^{\mathrm{K}}\varphi\) for every function \(\varphi \in \mathcal{H}(\mathrm{G},\mathrm{K})\) .

\begin{enumerate}
\def\labelenumi{\alph{enumi})}
\setcounter{enumi}{2}
\item
  Let \(C \subset G\) be a compact subset. There exists a function \(\varphi \in \mathcal{H}(G, K)\) such that \(\varphi(x) = 1\) for all \(x \in C\) .
\item
  Let (G,K) be a Gelfand pair. The group G is unimodular. (Note that the modulus of G belongs to \(\mathcal{H}(\mathrm{G},\mathrm{K})\) and interpret the integral of \(\varphi \in \mathcal{H}(\mathrm{G},\mathrm{K})\) as the value at \(e\) of a suitably chosen convolution.)
\item
  Suppose that there exists a continuous involution \(\theta\colon\mathrm{G}\to\mathrm{G}\) such that \(\mathrm{K}\theta(x)\mathrm{K}=\mathrm{K}x^{-1}\mathrm{K}\) for all \(x\in\mathrm{G}\) . Then (G,K) is a Gelfand pair. (Verify that the map \(f\mapsto\nu(f\circ\theta)\) on \(\mathcal{K}(\mathrm{G})\) is a left Haar measure on G, then that it is equal to \(\underline{\nu}\) ; deduce that \((f*g)\circ\theta=(f\circ\theta)*(g\circ\theta)\) and compare with the formula \(f*g=\check{g}*f\) .)
\item
  If G is compact, then \((\mathbf{G} \times \mathbf{G}, \Delta)\) is a Gelfand pair, where \(\Delta\) is the diagonal subgroup of \(G \times G\) image of the homomorphism \(x \mapsto (x, x)\) from G into \(G \times G\) .
\end{enumerate}

\begin{enumerate}
\def\labelenumi{\arabic{enumi})}
\setcounter{enumi}{25}
\tightlist
\item
  We resume the notation of the preceding exercise. We denote by \(\varphi_{\mathrm{K}}\) the characteristic function of K.
\end{enumerate}

\begin{enumerate}
\def\labelenumi{\alph{enumi})}
\item
  Let \(\pi\) be a unitary representation of G in a complex Hilbert space E. The representation of \(\mathcal{M}_c(\mathrm{G})\) in E defines, by passage to subspaces, a representation of the algebra \(\mathcal{H}(\mathrm{G},\mathrm{K})\) in the space \(\mathrm{E}^{\mathrm{K}}\) of K-invariant vectors of E.
\item
  If E admits a cyclic vector x such that \(x \in E^{K}\) , then \(\pi\) is irreducible.
\item
  Suppose that \(\pi\) is irreducible. If the space \(E^{K}\) is not reduced to 0, the representation of \(\mathcal{H}(G,K)\) in \(E^{K}\) is irreducible.
\end{enumerate}

\begin{enumerate}
\def\labelenumi{\alph{enumi})}
\setcounter{enumi}{3}
\tightlist
\item
  Let \(\varphi\) be a positive-definite function on G and let \((\varrho, x)\) be a cyclic Hilbertian realization of \(\varphi\) . We have \(\varphi \in \mathcal{H}(\mathrm{G}, \mathrm{K})\) if and only if \(x \in \mathrm{E}^{\mathrm{K}}\) .
\end{enumerate}

\begin{enumerate}
\def\labelenumi{\arabic{enumi})}
\setcounter{enumi}{26}
\tightlist
\item
  We resume the notation of the preceding exercise. We assume that (G, K) is a Gelfand pair.
\end{enumerate}

We say that \(\varphi\in\mathcal{H}(\mathrm{G},\mathrm{K})\) is a K-spherical function if the map

\[
f \mapsto \int_ {\mathrm{G}} f \check {\varphi} d \nu
\]

is a non-zero character of the algebra \(\mathcal{H}(\mathrm{G},\mathrm{K})\) .

\begin{enumerate}
\def\labelenumi{\alph{enumi})}
\tightlist
\item
  Let \(\varphi \in \mathcal{H}(\mathrm{G},\mathrm{K})\) . Prove that the following conditions are equivalent:
\end{enumerate}

\begin{enumerate}
\def\labelenumi{(\roman{enumi})}
\item
  The function \(\varphi\) is a K-spherical function;
\item
  For all \(x\) and \(y\) in \(\mathrm{G}\) , we have
\end{enumerate}

\[
\int_ {\mathrm{K}} \varphi (x k y) d \mu (k) = \varphi (x) \varphi (y);
\]

\begin{enumerate}
\def\labelenumi{(\roman{enumi})}
\setcounter{enumi}{2}
\tightlist
\item
  One has \(\varphi(e)=1\) and for every function \(\psi\in\mathcal{H}(\mathrm{G},\mathrm{K})\) , there exists \(c\in\mathbf{C}\) such that \(\psi*\varphi=c\varphi\) .
\end{enumerate}

\begin{enumerate}
\def\labelenumi{\alph{enumi})}
\setcounter{enumi}{1}
\item
  We denote by \(\mathcal{H}^1 (\mathrm{G},\mathrm{K})\) the space of \(f\in \mathrm{L}^1 (\mathrm{G})\) such that \(\varrho_{\mathrm{G}}(k_1,k_2)f = f\) for all \((k_{1},k_{2})\in \mathrm{K}\times \mathrm{K}\) . This is a Banach subalgebra of the involutive Banach algebra \(\mathrm{L}^1 (\mathrm{G})\) ; it is commutative.
\item
  The algebra \(\mathcal{H}(\mathrm{G},\mathrm{K})\) is identified with a dense subalgebra of \(\mathcal{H}^1 (\mathrm{G},\mathrm{K})\) .\\
\item
  The map which to \(\varphi \in \mathcal{H}(\mathrm{G},\mathrm{K})\) associates the linear form
\end{enumerate}

\[
f \mapsto \int_ {\mathrm{G}} f \check {\varphi} d \nu
\]

on \(\mathcal{H}^{1}(\mathrm{G},\mathrm{K})\) defines a bijection from the space of bounded K-spherical functions on G onto the space of non-zero characters of \(\mathcal{H}^{1}(\mathrm{G},\mathrm{K})\) . (Note that every character of \(\mathcal{H}^{1}(\mathrm{G},\mathrm{K})\) is continuous and of norm \(\leqslant 1\) , cf.~prop. 2 of I, p.~104.)

\begin{enumerate}
\def\labelenumi{\arabic{enumi})}
\setcounter{enumi}{27}
\tightlist
\item
  Let G be a locally compact group and K a compact subgroup of G.
\end{enumerate}

\begin{enumerate}
\def\labelenumi{\alph{enumi})}
\item
  The pair (G, K) is a Gelfand pair if and only if, for every unitary representation \(\pi\) of G in a complex Hilbert space E, the dimension of \(E^{K}\) is at most 1. (To prove that the condition is necessary, prove that an irreducible representation of a commutative involutive Banach algebra in a Hilbert space is of dimension 1; to prove that it is sufficient, use the cor. of th. 4 of V, p.~454 to verify that the characters of the algebra \(\mathcal{H}(G,K)\) separate its points.)
\item
  We assume that (G, K) is a Gelfand pair. Let \(\varphi \in \mathcal{H}(\mathrm{G},\mathrm{K})\) . Assume that \(\varphi\) is a positive-definite function and that \(\varphi(e) = 1\) . Show that \(\varphi\) is a K-spherical function if and only if \(\varphi\) has a Hilbert representation \((\varrho, x)\) such that \(\varrho\) is irreducible.
\item
  Suppose that \((\mathrm{G},\mathrm{K})\) is a Gelfand pair. The set of positive-definite K-spherical functions on G is in bijection with the set of irreducible representations \(\pi\in\widehat{G}\) such that \(E_{\pi}^{K}\) is nonzero.
\end{enumerate}

\begin{enumerate}
\def\labelenumi{\arabic{enumi})}
\setcounter{enumi}{28}
\tightlist
\item
  Let G be a compact group and K a closed subgroup of G. Suppose that (G, K) is a Gelfand pair. We denote by \(\nu\) (resp. \(\mu\) ) the normalized Haar measure on G (resp. on K).
\end{enumerate}

Let \(\varphi\) be a K-spherical function on G.

\begin{enumerate}
\def\labelenumi{\alph{enumi})}
\item
  We denote by \(E_{\varphi}\) the subspace of \(\mathcal{C}(G)\) functions of the form \(\alpha * \varphi\) , where \(\alpha\) is a measure on G with finite support. For every \(\psi \in \mathcal{H}(G, K)\) , the endomorphism \(v_{\psi}: f \mapsto f * \psi\) of \(\mathcal{C}(G)\) is compact and \(E_{\varphi}\) is contained in an eigenspace of \(v_{\psi}\) .
\item
  The space \(E_{\varphi}\) is of finite dimension. It identifies with a subspace of G-finite vectors of the regular representation \(\gamma_{G}\) containing \(\varphi\) .
\item
  The space of K-invariant vectors of \(E_{\varphi}\) is of dimension 1 and the representation of G in \(E_{\varphi}\) is irreducible.
\end{enumerate}

\begin{enumerate}
\def\labelenumi{\alph{enumi})}
\setcounter{enumi}{3}
\tightlist
\item
  Let \((f_i)_{i\in I}\) be an orthonormal basis of \(\mathrm{E}_{\varphi}\) and \(k\colon \mathbf{G}\times \mathbf{G}\to \mathbf{C}\) the function defined by

\[
k (x, y) = \sum_ {i \in \mathrm{I}} \overline {{f _ {i} (x)}} f _ {i} (y).
\]

One has

\[
f (x) = \int_ {\mathrm{G}} k (x, y) f (y) d \nu (y)
\]

for all \(f \in \mathrm{E}_{\varphi}\) and \(x \in \mathrm{G}\) .

\end{enumerate}

\begin{enumerate}
\def\labelenumi{\alph{enumi})}
\setcounter{enumi}{4}
\item
  One has \(k(x,y)=\dim(\mathrm{E}_{\varphi})\varphi(y^{-1}x)\) for all \((x,y)\in\mathrm{G}\times\mathrm{G}\) .
\item
  The function \(\varphi\) is of positive type.
\item
  Let \(\varphi\) be a K-spherical function on G and let \((\pi, x)\) be a Hilbert realization of \(\varphi\) . The representation \(\pi\) is irreducible; we have
\end{enumerate}

\[
\varphi (x) = \int_ {\mathrm{K}} \chi_ {\pi} (x ^ {- 1} k) d \mu (k)
\]

for all \(x \in \mathrm{G}\) , and \(\int_{\mathrm{G}} |\varphi|^2 d\nu = \dim(\pi)^{-1}\) .

\begin{enumerate}
\def\labelenumi{\arabic{enumi})}
\setcounter{enumi}{29}
\tightlist
\item
  Let \(n \in N\) . Let us denote by G the subgroup of the affine group \(\mathrm{A}(\mathbf{R}^{n})\) of the space \(R^{n}\) formed by the elements \(g \in \mathrm{A}(\mathbf{R}^{n})\) such that the linear map \(g_{0}\) associated with g belongs to \(\mathbf{SO}(\mathbf{R}^{n})\) ; it is a real Lie group (cf. A, II, p. 131 and LIE, III, p. 102, example). We denote by \(g = (g_{0}, a)\) the elements of G, with \(g_{0} \in \mathbf{SO}(\mathbf{R}^{n})\) and \(a \in R^{n}\) .
\end{enumerate}

\begin{enumerate}
\def\labelenumi{\alph{enumi})}
\item
  The pair \((\mathbf{G},\mathbf{SO}(\mathbf{R}^n))\) is a Gelfand pair.
\item
  Let \(\varphi \in \mathcal{H}(\mathrm{G},\mathbf{SO}(\mathbf{R}^n))\) . Denote by \(f\colon \mathbf{R}^n\to \mathbf{C}\) the function defined by \(f(r) = \varphi (1_{\mathbf{R}^N},g(e_1))\) . The map \(u\colon \varphi \mapsto f\) is an injective morphism of complex algebras from \(\mathcal{H}(\mathrm{G},\mathrm{K})\) into \(\mathcal{K}(\mathbf{R}^n)\) ; its image is the subspace of \(\mathcal{K}(\mathbf{R}^n)\) consisting of functions \(f\) such that \(f(g_0(x)) = f(x)\) for all \(g_0\in \mathbf{SO}(\mathbf{R}^n)\) and every \(x\in \mathbf{R}^n\) .
\item
  We denote by \(\Delta\) the Laplacian on \(\mathbf{R}^n\) . A function \(\varphi \in \mathcal{K}(\mathrm{G})\) is a spherical function if and only if the function \(f = u(\varphi)\) is of class \(\mathrm{C}^{\infty}\) , satisfies \(f(0) = 1\) , and if there exists \(\lambda \in \mathbf{C}\) such that \(\Delta f = \lambda f\) (Use the fact that \(\Delta(f \circ g_0) = (\Delta f) \circ g_0\) for all \(g_0 \in \mathbf{SO}(\mathbf{R}^n)\) .)
\end{enumerate}

\begin{enumerate}
\def\labelenumi{\arabic{enumi})}
\setcounter{enumi}{30}
\tightlist
\item
  Let \(n\) be an integer \(\geqslant 1\) . We denote by N the Euclidean norm on \(\mathbf{R}^n\) , and we denote by \(\lambda_n\) the Lebesgue measure on \(\mathbf{R}^n\) . We denote by \(\omega_n\) the unique positive measure of total mass 1 on the sphere \(\mathbf{S}_n \subset \mathbf{R}^{n+1}\) invariant under the action of the orthogonal group \(\mathbf{O}(n+1, \mathbf{R})\) (cf.~INT, VIII, p.~116, exercise 8).
\end{enumerate}

We identify \(R^{n}\) with the complement of a projective hyperplane H in the projective space \(\mathbf{P}_{n}(\mathbf{R})\) (cf.~TG, VI, p.~15, prop. 5), and we shall interpret the measures on \(R^{n}\) as measures on \(\mathbf{P}_{n}(\mathbf{R})\) supported on the complement of H.

For every \(g \in \mathbf{GL}(n+1, \mathbf{R})\) , we denote by \(\widetilde{g}\) the homeomorphism of \(\mathbf{P}_{n}(\mathbf{R})\) defined by g.

\begin{enumerate}
\def\labelenumi{\alph{enumi})}
\tightlist
\item
  The positive measure \(c_{n}\) on \(\mathbf{R}^n\) defined by
\end{enumerate}

\[
c _ {n} = \pi^ {- (n + 1) / 2} \Gamma \Big (\frac {n + 1}{2} \Big) (1 + \mathrm{N} ^ {2}) ^ {- (n + 1) / 2} \lambda_ {n}
\]

(“Cauchy measure”) is bounded and has total mass 1. If \(n = 1\) this is the stable measure of parameter 1 (cf.~exercise V, p.~496). For every projective hyperplane \(\mathrm{H}'\) of \(\mathbf{P}_n(\mathbf{R})\) , we have \(c_n(\mathrm{H}') = 0\) .

\begin{enumerate}
\def\labelenumi{\alph{enumi})}
\setcounter{enumi}{1}
\item
  The canonical map from \(\mathbf{R}^{n+1} - \{0\}\) into \(\mathbf{P}_n(\mathbf{R})\) defines, by passing to subspaces, a continuous surjective map \(\sigma: \mathbf{S}_n - (\mathbf{R}^n \times \{0\}) \to \mathbf{R}^n\) . We have \(\sigma(\omega_n) = c_n\) .
\item
  Let \(g \in \mathbf{GL}(n+1, \mathbf{R})\) . There exists \(a \in \mathbf{GL}(n, \mathbf{R})\) and \(b \in R^n\) such that the image measure \(\widetilde{g}(c_n)\) is equal to the image of \(c_n\) by the affine map \(x \mapsto a(x) + b\) .
\end{enumerate}

In what follows, we propose to prove a converse of this assertion.

\begin{enumerate}
\def\labelenumi{\alph{enumi})}
\setcounter{enumi}{3}
\tightlist
\item
  Let \(\mu\) be a bounded positive measure on \(\mathbf{P}_{n}(\mathbf{R})\) such that every projective hyperplane is \(\mu\) -negligible. The set of \(g \in \mathbf{GL}(n+1, \mathbf{R})\) such that \(g(\mu) = \mu\) and \(|\det(g)| = 1\) is a compact subgroup of \(\mathbf{GL}(n+1, \mathbf{R})\) .

\end{enumerate}

\begin{enumerate}
\def\labelenumi{\alph{enumi})}
\setcounter{enumi}{4}
\item
  Let \(\mu\) be a bounded positive measure on \(\mathbf{P}_n(\mathbf{R})\) of total mass 1 such that, for all \(g \in \mathbf{GL}(n + 1, \mathbf{R})\) there exists an affine map \(f\) of \(\mathbf{R}^n\) such that \(g(\mu) = f(\mu)\) . Every projective hyperplane is \(\mu\) -negligible.
\item
  Let \(K_{\mu}\) be the subgroup of \(g \in \mathbf{GL}(n+1, \mathbf{R})\) such that \(g(\mu) = \mu\) . There exists \(h \in \mathbf{GL}(n+1, \mathbf{R})\) such that \(K_{\mu} \subset h \mathbf{O}(n+1, \mathbf{R}) h^{-1}\) .
\item
  Let \(K = h^{-1}K_{\mu}h\) . We have \(\mathbf{O}(n + 1, \mathbf{R}) = \mathbf{G} \cdot \mathbf{K}\) , where G is the subgroup of \(\mathbf{O}(n + 1, \mathbf{R})\) consisting of the matrices
\end{enumerate}

\[
\left( \begin{array}{c c} x & 0 \\ 0 & 1 \end{array} \right)
\]

with \(x\in \mathbf{O}(n,\mathbf{R})\)

\begin{enumerate}
\def\labelenumi{\alph{enumi})}
\setcounter{enumi}{7}
\tightlist
\item
  The measure \(h^{-1}(\mu)\) coincides with \(c_{n}\) . (Observe that the group K acts transitively on \(S_{n}\) and that the image measure \(\sigma(h^{-1}(\mu))\) is a K-invariant measure on \(S_{n}\) .)
\end{enumerate}

\begin{enumerate}
\def\labelenumi{\arabic{enumi})}
\setcounter{enumi}{31}
\tightlist
\item
  Let G be a locally compact group. Suppose there exists a subgroup H of G that is compact and open in G. We denote by \(\mu\) the unique left Haar measure on G such that \(\mu(\mathrm{H}) = 1\) and \(\varphi\) the characteristic function of H.
\end{enumerate}

\begin{enumerate}
\def\labelenumi{\alph{enumi})}
\item
  Let \(\pi\) be a square-integrable representation of G in a complex Hilbert space E. Show that the endomorphism \(\pi(\varphi)\) is a Hilbert-Schmidt endomorphism.
\item
  The space \(E^{H}\) is finite-dimensional and satisfies
\end{enumerate}

\[
\dim (\mathrm{E} ^ {\mathrm{H}}) \leqslant \frac {1}{c (\pi)}
\]

where \(c(\pi)\) is the formal degree of \(\pi\) relative to \(\{e\}\) (Use Exercise 6 of V, p.~487.)

\begin{enumerate}
\def\labelenumi{\alph{enumi})}
\setcounter{enumi}{2}
\tightlist
\item
  A discrete infinite group has no square-integrable representation; if \(n \geqslant 2\) is an integer, the group \(\mathbf{SL}(n, \mathbf{Z})\) has no square-integrable representation modulo the center.
\end{enumerate}

\specialchapter{HISTORICAL NOTE}

(Chapters I to V)

The notions of eigenvalue and Fourier series were already well known in 1830, permeating many fields of analysis throughout the nineteenth century (cf.~ÉHM, pp.~114, 260, 275). But spectral theory and harmonic analysis, in the sense in which we understand them in this book, did not begin to take their present form until about a century later --- at the same time as their junction with the study of group representations around 1930, and then with that of normed algebras around 1940.

We confine ourselves here to sketching the main lines of their evolution in the period from 1906—when Hilbert introduced the notion of the “continuous spectrum” in his work on integral equations—to the years 1945–1950, when the theories presented here essentially acquired the form they have retained to the present day.

\histitem{Discovery of the continuous spectrum}

In the Note on topological vector spaces (cf. ÉHM, pp. 262–263), we described how I. Fredholm, around 1900 — continuing the work of C. Neumann, V. Volterra, and H. Poincaré — came to consider the equation

\[
u (x) - \lambda \int_ {\mathrm{I}} \mathrm{K} (x, y) u (y) d y = f (x)\tag{1}
\]

with unknown function \(u\colon I\to\mathbf{C}\) , in the case where the interval \(I\subset\mathbf{R}\) is compact and where the kernel \(K\) is continuous on \(I\times I\) . A skillful use of "infinite determinants," based on ideas of Poincaré and H. von Koch, allowed him to obtain families of solutions whose dependence on the complex variable \(\lambda\) is meromorphic, then to prove the famous “alternative” regarding the existence of solutions (cf.~III, p.~81, th. 5). But Fredholm immediately specializes his results to the case \(\lambda = -1\) without exploiting the fact that equation (1) is an eigenvalue problem \(^{(1)}\) .

When the kernel K is symmetric, we also saw how Hilbert, in his first memoir on integral equations [31], recognized the kinship between the Fredholm problem and the search for the “principal axes” of a quadratic form. Starting from this reduction for finite “sections” (discretizations) of the kernel K, he obtained by passage to the limit the formula

\[
\int_ {\mathrm{I}} \mathrm{K} (s, t) x (s) x (t) d t = \sum_ {n = 1} ^ {\infty} \frac {1}{\lambda_ {n}} \int_ {\mathrm{I}} \varphi_ {n} (s) x (s) d s,\tag{2}
\]

where the \(\lambda_{n}\) are the eigenvalues of the kernel K, the \(\varphi_{n}\) form an orthonormal system of associated eigenvectors \(^{(2)}\) and where the equality holds as soon as x is square-integrable on I (cf.~ÉHM, p.~264).

The principal obstacles facing Hilbert concerned the passage from finite sums to the integrals of formula (2). He deduced from this formula the existence of the eigenvalues — however, he gave for the latter a variational characterization independent of the passage to the limit, inspired by the classical method for determining eigenvalues in finite dimension (cf.~no. 4 of IV, p.~153). E. Schmidt would show in 1905, drawing on the work of H. Schwarz, how to obtain (2) without going through the “infinite determinants” on which the methods of Fredholm and Hilbert depended (cf.~ÉHM, p.~265).

But already in 1906, Hilbert [32] turns to the more general problem of reducing a quadratic form

\[
\mathrm{B} (x, x) = \sum_ {p, q = 0} ^ {\infty} b _ {p q} x _ {p} x _ {q}, \qquad x = (x _ {p}) _ {p \in \mathbf {N}}
\]

on \(E = \ell_{\mathbf{C}}^{2}(\mathbf{N})\) . He certainly has in mind that by passing to coefficients in an orthonormal basis of E, problem (1) reduces to the search for solutions of the “infinite linear system”

\[
x _ {p} + \sum_ {q = 0} ^ {\infty} b _ {p q} x _ {q} = f _ {p} \quad (p \in \mathbf {N}).
\]

He notes, however, that for the Fredholm problem one deals with very particular quadratic forms, for which \(\mathrm{B}(x_{n}, x_{n})\) tends to \(\mathrm{B}(x, x)\) as soon as \(x_{n}\) converges weakly to x; he calls them “completely continuous” (vollstetig).

Hilbert insists on the essential difference that distinguishes them from quadratic forms bounded on E, and decides to undertake the more general reduction of the latter. Again, his method is to start from the reduction of the "finite sections."

\[
(x _ {0}, \ldots , x _ {n}) \mapsto \sum_ {p, q = 0} ^ {n} b _ {p q} x _ {p} x _ {q}
\]

of the form B, and then to let n tend to infinity. For this passage to the limit, he takes advantage of the ideas on integration introduced by T. Stieltjes in his work on continued fractions [71]. He shows that every bounded quadratic form can, after an orthogonal transformation of the variables, be put in the form

\[
\sum_ {p \in {\bf N}} \frac {1}{\lambda_ {p}} x _ {p} ^ {2} + \int_ {s \in {\bf R}} \frac {1}{s} d \sigma (s, x).\tag{3}
\]

In this formula, the \(\lambda_{p}\) are the eigenvalues of B and \(x\mapsto\sigma(s,x)\) is (for a fixed \(s\in\mathbf{R}\)) a positive definite separating quadratic form on E, while \(s\mapsto\sigma(s,x)\) is \(^{(3)}\) (for a fixed \(x\in\mathrm{E}\)) a continuous increasing function that tends to 0 at \(-\infty\) and toward \(\|x\|^2\) at \(+\infty\) (cf. ÉHM, p. 284, for Stieltjes notation).

Let S denote the set of points of \(\overline{R}\) that do not admit a neighborhood in which \(s \mapsto \sigma(s, x)\) remains constant for all x; then the integral in (3) can naturally be taken over S. Hilbert calls the set S the “continuous spectrum” (Streckenspektrum), the set of eigenvalues of B the “point spectrum” (Punktspektrum), and their union the “spectrum” (Spektrum). This term, coming from optics, had been used in 1897 by W. Wirtinger [91] in the study of differential equations with periodic coefficients.

Hilbert would quickly draw from his “spectral” methods a profound renewal of several problems of classical analysis, pertaining to the study of ordinary differential equations (in particular of Sturm–Liouville type), partial differential equations, or functions of a complex variable. We will not report here on the profusion of results that followed in the first quarter of the twentieth century and refer to the synthesis by E. Hellinger and O. Toeplitz [29]. We shall nevertheless mention some of these works on integral equations, in which the germs of later developments of the abstract theory are to be found.

For Hilbert and his students, faithful to a German tradition in linear algebra, the model remains the principal axes theorem; thus it is always a matter of studying the spectra of bilinear or quadratic forms, in particular on the “Hilbert” space \(E = \ell_{\mathbf{C}}^{2}(\mathbf{N})\) . The fact that one can associate any continuous bilinear form B on E with the continuous endomorphism A of E characterized by \(\langle Ax, y \rangle = B(x, y)\) was known to the Hilbert school, particularly after the works of E. Schmidt, M. Fréchet, and F. Riesz in 1907–1908. But it was F. Riesz who showed, in an admirable book [59] from 1913, the interest of placing endomorphisms at the forefront in this context.

Riesz gives the modern definition of the spectrum of an element A of \(\mathcal{L}(\mathrm{E})\) [59, p.~139], notes that this is a closed, bounded set, and shows that the resolvent \(\lambda \mapsto (\lambda 1_{\mathrm{E}} - \mathrm{A})^{-1}\) is holomorphic on the complement of the spectrum of A. Above all, he assigns a central role to the algebra structure of \(\mathcal{L}(\mathrm{E})\): Drawing on the work of Volterra, he developed a first version of the functional calculus for symmetric endomorphisms, which allowed him to recover simply the reduction results of Hilbert and Schmidt. If A is a symmetric element of \(\mathcal{L}(\mathrm{E})\) , Riesz shows how to define \(f(\mathrm{A})\) for every function \(f\) semicontinuous (from below or from above) on \(\mathbf{R}\) in \(\mathbf{R}\) , and notes that \(f \mapsto f(\mathrm{A})\) is what we today call an algebra morphism [59, p.~129--130]. Applying this idea to the characteristic function of an interval \([\xi, +\infty[\) , one obtains for every \(\xi\) a symmetric endomorphism \(\mathrm{A}_{\xi}\) of E; using the Stieltjes integral, Riesz then proves a version of Hilbert's relation (3):

\[
\langle \mathrm{A} x, y \rangle = \int_ {\mathbf {R}} \xi d \langle \mathrm{A} _ {\xi} x, y \rangle
\]

for \(x, y \in \mathrm{E}\) , an equality that he even writes

\[
\mathrm{A} = \int_ {\mathbf {R}} \xi d \mathrm{A} _ {\xi}.
\]

The spectrum of A is then the set of values \(\xi_{0}\) such that \(A_{\xi}\) does not remain constant in any neighborhood of \(\xi_{0}\) and the point spectrum that of the points of discontinuity of \(\xi \mapsto A_{\xi}\) .

Shortly before concluding with the applications of the theory, Riesz notes [59, p.~146] that its results take a particularly simple form when the endomorphism A arises from one of Hilbert's "completely continuous" bilinear forms: the spectrum reduces to the set of eigenvalues (to which 0 must possibly be added), the eigenvalues are arranged in a sequence tending to 0, and each nonzero eigenvalue has finite multiplicity (III, p.~90, prop. 5).

\histitem{Compact operators}

Riesz considerably amplifies these remarks less than three years later, in his memoir on functional equations [60] which develops, with a clarity undimmed after more than a century, the essentials of the spectral theory of compact operators on Banach spaces. The stated aim of this text is to reformulate Fredholm theory for equation (1) on a compact interval I, by considering the space \(E = \mathcal{C}(I)\) equipped with the topology of uniform convergence. But as we have already indicated in an earlier note (ÉHM, p.~268), Riesz is clearly aware that his results apply to any Banach space — a notion that would only be introduced in the following decade.

Inspired by Fréchet's ideas on general topology, Riesz now considers the endomorphisms of E that map every bounded sequence to a relatively compact sequence. Like Hilbert, he still calls them "completely continuous" and notes that this new notion is equivalent, in the case of \(\mathrm{E} = \ell_{\mathbf{C}}^{2}(\mathbf{N})\) , to the notion of complete continuity that he had used in his 1913 book \(^{(4)}\) .

Riesz studies endomorphisms of the form \(B = 1_{E} - A\) where A is a completely continuous endomorphism of E, and derives all their spectral properties from the fact that a locally compact normed vector space must be finite-dimensional — a fact apparently discovered for the occasion. He proves easily that the kernel of B is finite-dimensional and that its image is closed and of finite codimension. As E. Weyr [87] had done it in finite dimension, he then considers the iterates \(B^{k}\) . He shows that the sequence of their kernels and that of their images are stationary, then introduces what we today call the nilspace and conilspace of B, efficiently showing that E is their topological direct sum. Applying this observation to endomorphisms \(B_{\lambda} = 1_{E} - \lambda A\) for \(\lambda \in C\) , he deduces without difficulty the spectral properties of completely continuous endomorphisms of a Banach space, and this in a more or less definitive form.

Among the main results on the spectrum of these operators, only those that would require a more mature view of the notion of function space, particularly of duality, seem to escape Riesz. For example, in the late 1920s, T. Hildebrandt [33] and J. Schauder [63] will show (within the now well-established framework of Banach spaces) that the adjoint of a completely continuous endomorphism is completely continuous.

Moreover, the notion of Riesz complete continuity still depends on the use of sequences; Banach [4] will free himself from it in his 1932 book to consider mappings that transform every bounded set into a relatively compact set. The term “compact” linear mapping would only gradually establish itself in the second half of the 20th century, presumably following E. Hille's book [34] on semigroups of operators.

Several questions raised in the 1930s would remain open for a long time. In this direction, we note the resolution [18] in 1973 of the "approximation problem" made famous by Banach and S. Mazur in 1932 and 1938 (see Remark 6 of III, p.~16 and Exercise 25 of III, p.~112).

At the same time, Lomonosov [40] proves the existence of nontrivial closed invariant subspaces for an endomorphism commuting with a nonzero compact endomorphism in a locally convex space (cor. 2 of III, p.~13). This result was one of the most significant advances on the problem of the existence of such a subspace for an arbitrary continuous endomorphism (the “invariant subspace problem”). This problem is still open at present in the case of continuous endomorphisms of separable Hilbert spaces; P. Enflo [19] constructed in 1987 the first example of an endomorphism of a Banach space admitting no nontrivial closed invariant subspace.